% =============================================================
% Capítulo — Desenvolvimento
% =============================================================
% Reescrito em 22/08/2026: o conteúdo anterior descrevia o experimento
% com ResNet-50 / CIFAR-100, anterior ao pivô para Transformer decidido
% na reunião de 06/05/2026. Esta versão descreve o experimento
% efetivamente implementado, com GPT-2 small sobre o WikiText-2.
% =============================================================

\chapter{Desenvolvimento}\label{chap:desenvolvimento}

% AJUSTE: revisão 30/09 — sai a frase de abertura que retomava o capítulo anterior; a
% referência aos critérios passa para a frase seguinte, que ficaria sem antecedente.
% AJUSTE: revisão 30/09 — o capítulo passa a seguir as mesmas quatro etapas da
% Metodologia, na mesma ordem e com os mesmos nomes (esboço do orientador). Base de dados
% e métricas entram na Etapa 1; parâmetros podáveis, na Etapa 2. A antiga subseção "As
% Quatro Etapas Centrais" foi distribuída pelas aberturas das etapas.
Este capítulo descreve as decisões concretas que decorrem dos critérios estabelecidos no
Capítulo~\ref{chap:metodologia} e o experimento tal como foi efetivamente implementado.
Após uma visão geral (Seção~\ref{sec:visao_geral}), o capítulo segue as mesmas quatro
etapas da metodologia, na mesma ordem: Etapa 1 --- Modelo Base
(Seção~\ref{sec:etapa1}); Etapa 2 --- Poda sem Otimização (Seção~\ref{sec:metodos_poda});
Etapa 3 --- Poda com Otimização (Seção~\ref{sec:estrategia_otimizada}); e Etapa 4 ---
\textit{Fine-tuning} de Recuperação (Seção~\ref{sec:finetuning}). As seções finais tratam
de aspectos comuns a todas as etapas: o \textit{framework} e as bibliotecas empregadas, a
organização do \textit{pipeline} experimental, o ambiente de execução e as limitações
reconhecidas.

\section{Visão Geral}\label{sec:visao_geral}

A pergunta que orienta o trabalho é empírica e comparativa. Não se pretende demonstrar que
a poda comprime modelos, resultado já consolidado pela literatura
\cite{Han2015, Li2017, Blalock2020}, mas responder a duas questões mais específicas,
sob condições controladas e reprodutíveis:

\begin{enumerate}
    \item[(A)] Em que medida a redução \textit{teórica} de parâmetros se converte em
          redução \textit{real} de tempo de inferência e de consumo de energia, e como
          essa conversão depende da granularidade da poda?
    \item[(B)] Uma alocação não uniforme da poda entre as camadas, definida por uma análise
          de sensibilidade por camada, melhora o compromisso entre compressão e qualidade em
          relação às alocações sem análise de sensibilidade?
\end{enumerate}

Responder a essas duas perguntas exige que modelos comprimidos por estratégias distintas
sejam avaliados exatamente pelo mesmo conjunto de métricas, na mesma base de dados e no mesmo
\textit{hardware}. Essa simetria de avaliação é condição necessária para que qualquer
comparação seja interpretável, e é o princípio que organiza todo o desenvolvimento
descrito adiante.

\subsection{Discussão das Decisões Metodológicas}\label{subsec:discussao_decisoes}

Todo desenho experimental implica escolhas que condicionam o escopo dos resultados. A
Tabela~\ref{tab:decisoes_metodologicas} sintetiza as principais decisões adotadas, as
alternativas consideradas e a justificativa de cada uma. As seções seguintes desenvolvem
em maior profundidade as que são mais críticas para a interpretação dos resultados.

\begin{table}[htbp]
\centering
\caption{Síntese das principais decisões metodológicas do trabalho.}
\label{tab:decisoes_metodologicas}
\small
\begin{tabular}{>{\raggedright\arraybackslash}p{2.4cm} >{\raggedright\arraybackslash}p{3.4cm} >{\raggedright\arraybackslash}p{2.6cm} >{\raggedright\arraybackslash}p{5.1cm}}
\hline
\textbf{Decisão} & \textbf{Alternativas consideradas} & \textbf{Escolha} & \textbf{Justificativa} \\
\hline

Base de dados
& WikiText-103, Penn Treebank, OpenWebText, The Pile
& \textbf{WikiText-2}
& Atende aos quatro critérios da Seção~\ref{subsec:metodologia_base}; entre os que
  atendem, tem o menor conjunto de treino (Tabela~\ref{tab:candidatos_dataset}). \\

Modelo
& GPT-2 \textit{medium}/\textit{large}/XL, DistilGPT2, OPT-125M, Pythia-160M, BERT, Llama 2
& \textbf{GPT-2 \textit{small}}
& Atende aos seis critérios da Seção~\ref{subsec:metodologia_selecao_modelo} e forma com o
  WikiText-2 um par que atende a D1--D4; entre os que atendem, é o de menor número de
  parâmetros (Tabela~\ref{tab:candidatos_modelo}). \\

Estratégias de poda
& SparseGPT, Wanda, \textit{movement pruning}, poda de cabeças por gradiente
& \textbf{Duas} (magnitude e estruturada)
& Únicas alternativas que atendem a P1 e P2; juntas, atendem a P3
  (Tabela~\ref{tab:candidatos_poda}). A restrição de tempo experimental inviabilizaria uma
  terceira com a mesma profundidade. \\

Regime de poda
& Iterativo (podar e retreinar em ciclos)
& \textbf{\textit{One-shot} primeiro}
& Isola o efeito da poda do efeito do retreinamento; o \textit{fine-tuning} entra depois,
  como etapa separada e mensurável. \\

Estratégia otimizada
& Taxas dinâmicas, variação do momento da poda, destilação de conhecimento, quantização
& \textbf{Sensibilidade por camada}
& Única que atende a O1--O3 (Seção~\ref{sec:metodologia_poda_com}): altera apenas a
  distribuição da poda entre as camadas, objeto da hipótese do trabalho. \\

Métrica principal
& \textit{Accuracy} em tarefas descendentes
& \textbf{Perplexity}
& Métrica intrínseca padrão para modelos de linguagem, mensurável sem
  \textit{fine-tuning} em tarefa específica. \\

Ambiente
& Máquina local, Google Colab
& \textbf{Kaggle (GPU T4)}
& Acesso gratuito a GPU dedicada com sessões longas o bastante para as varreduras
  completas. \\
\hline
\end{tabular}\\
\textit{Fonte: Elaboração própria.}
\end{table}

\section{Etapa 1 --- Modelo Base}\label{sec:etapa1}

Na primeira etapa, o GPT-2 \textit{small} \cite{Radford2019} foi selecionado como modelo
base. Utilizou-se a versão pré-treinada disponibilizada na plataforma Hugging
Face, executada sem nenhuma modificação e avaliada no conjunto de teste da base de dados
WikiText-2 \cite{Merity2017}. Registraram-se a \textit{perplexity}, a contagem de
parâmetros, o tamanho em memória, os FLOPs, o tempo de inferência e o consumo de energia:
as três primeiras métricas foram obtidas com a biblioteca PyTorch; os FLOPs, com a
biblioteca \texttt{thop}; o tempo de inferência, com o temporizador \texttt{perf\_counter}
da linguagem Python; e o consumo de energia, com a biblioteca CodeCarbon
(Tabela~\ref{tab:metricas}). O modelo base obteve \textit{perplexity} de 25,17
(Tabela~\ref{tab:gpt2_config}), valor contra o qual são comparadas todas as configurações
podadas das etapas seguintes. As subseções a seguir apresentam a seleção do modelo
(Seção~\ref{sec:modelo}), a seleção da base de dados (Seção~\ref{sec:dataset}) e as
métricas registradas (Seção~\ref{sec:metricas}).

\subsection{Seleção do Modelo}\label{sec:modelo}

\subsubsection{O GPT-2 \textit{small}}\label{subsec:gpt2small}

O modelo base é o GPT-2 \textit{small}, proposto por \citeonline{Radford2019} e
descrito na Seção~\ref{subsec:llm_gpt2}. Trata-se de um modelo autorregressivo composto
exclusivamente pelo componente decodificador da arquitetura \textit{Transformer}. Os pesos
pré-treinados foram obtidos pela plataforma Hugging Face, que os disponibiliza publicamente,
e carregados por meio da biblioteca PyTorch \cite{Paszke2019}.

% AJUSTE: 24/10 — "camada" e "bloco" alternavam para a mesma unidade (ex.: "primeira
% camada" e "bloco 0" no mesmo parágrafo dos Resultados), e a Figura 8 numerava de 1 a 12
% enquanto as tabelas usam 0 a 11. Define-se aqui, uma vez, que cada bloco é uma camada,
% numerada de 0 a 11 como no código; daqui em diante, "camada" para essa unidade.
A Tabela~\ref{tab:gpt2_config} resume a configuração da variante utilizada. Ao longo deste
trabalho, cada um dos 12 blocos \textit{Transformer} do modelo é denominado
\textbf{camada}, e as camadas são numeradas de 0 a 11, na ordem em que processam a
entrada, conforme a indexação adotada na implementação. É nessa unidade que se definem a
análise de sensibilidade por camada e as taxas de poda por camada.

\begin{table}[htb]
\centering
\caption{Configuração do GPT-2 \textit{small} utilizado como modelo base.}
\label{tab:gpt2_config}
\small
\begin{tabular}{p{6.5cm} r}
\hline
\textbf{Característica} & \textbf{Valor} \\
\hline
Blocos \textit{Transformer} empilhados (camadas) & 12 \\
Cabeças de atenção por camada & 12 \\
Dimensão do modelo ($d_{\text{model}}$) & 768 \\
Dimensão intermediária do MLP ($d_{\text{ff}}$) & 3.072 \\
Tamanho do vocabulário & 50.257 \\
Janela de contexto máxima & 1.024 \textit{tokens} \\
Total de parâmetros & 124.439.808 \\
Tamanho em memória & 474,70 MB \\
FLOPs por janela de 512 \textit{tokens} & 126,6 G \\
Perplexity no WikiText-2 (teste) & 25,17 \\
\hline
\end{tabular}\\
\textit{Fonte: Elaboração própria.}
\end{table}

\subsubsection{Justificativa da Escolha}\label{subsec:justificativa_gpt2}

A Seção~\ref{subsec:metodologia_selecao_modelo} estabeleceu seis critérios objetivos (C1 a C6) para a
seleção do modelo base. A Tabela~\ref{tab:candidatos_modelo} aplica esses
critérios a nove candidatos: as quatro variantes do GPT-2 \cite{Radford2019}, o
DistilGPT2 \cite{Sanh2019}, o OPT-125M \cite{Zhang2022}, o Pythia-160M
\cite{Biderman2023}, o BERT-base \cite{Devlin2019} e o Llama 2 7B \cite{Touvron2023}. Os
candidatos compreendem as quatro variantes da família GPT-2; três modelos abertos de porte
próximo ao do GPT-2 \textit{small}; e dois modelos de classes distintas, incluídos para
verificar os limites dos critérios: um codificador e um modelo com bilhões de parâmetros.
Como a seleção é conjunta (Seção~\ref{subsec:metodologia_selecao_modelo}), a tabela verifica também
o critério D4 de cada candidato em relação ao WikiText-2.
\todo[inline]{Avaliar a inclusão de modelos pequenos posteriores a 2023 (por exemplo,
SmolLM2, Qwen2.5-0.5B, Llama 3.2 1B), que a banca pode questionar.}

\begin{table}[htb]
\centering
\caption{Candidatos ao modelo base e verificação dos critérios de seleção.}
\label{tab:candidatos_modelo}
\footnotesize
\setlength{\tabcolsep}{4pt}
\begin{tabular}{l r r r c c c c c c c}
\hline
\textbf{Modelo} & \textbf{Parâm.} & \textbf{Mem.\ FT (GB)} & \textbf{Downl.\ (mi)} & \textbf{C1} & \textbf{C2} & \textbf{C3} & \textbf{C4} & \textbf{C5} & \textbf{C6} & \textbf{D4} \\
\hline
GPT-2 \textit{small}  & 124M  & 2,0   & 930,1   & \checkmark & \checkmark & \checkmark & \checkmark & \checkmark & \checkmark & \checkmark \\
GPT-2 \textit{medium} & 355M  & 5,7   & 29,8    & \checkmark & \checkmark & \checkmark & \checkmark & \checkmark & \checkmark & \checkmark \\
GPT-2 \textit{large}  & 774M  & 12,4  & 61,0    & \checkmark & \checkmark & \checkmark & \checkmark & \checkmark & \checkmark & \checkmark \\
GPT-2 XL              & 1,56B & 24,9  & 17,0    & \checkmark & \checkmark & $\times$   & \checkmark & \checkmark & \checkmark & \checkmark \\
DistilGPT2            & 82M   & 1,3   & 256,8   & \checkmark & \checkmark & \checkmark & $\times$   & ---        & \checkmark & ---        \\
OPT-125M              & 125M  & 2,0   & 215,8   & \checkmark & \checkmark & \checkmark & \checkmark & \checkmark & \checkmark & $\times$   \\
Pythia-160M           & 162M  & 2,6   & 27,2    & \checkmark & \checkmark & \checkmark & \checkmark & ---        & \checkmark & $\times$   \\
BERT-base             & 110M  & 1,8   & 3.246,2 & $\times$   & \checkmark & \checkmark & \checkmark & $\times$   & \checkmark & $\times$   \\
Llama 2 7B            & 6,74B & 107,8 & 34,6    & \checkmark & $\times$   & $\times$   & \checkmark & \checkmark & \checkmark & ---        \\
\hline
\end{tabular}
\par\smallskip
\begin{minipage}{0.92\textwidth}\centering C1: arquitetura decodificadora autorregressiva; C2: pesos abertos; C3: viabilidade
do \textit{fine-tuning}; C4: modelo original; C5: resultado de referência publicado; C6:
implementação disponível (critérios definidos na Seção~\ref{subsec:metodologia_selecao_modelo}); D4:
ausência de sobreposição entre o pré-treinamento e o WikiText-2
(Seção~\ref{subsec:metodologia_base}). ---: não verificado, candidato já eliminado por
outro critério. Mem.\ FT: memória mínima para o \textit{fine-tuning} completo, a 16 bytes
por parâmetro. Downl.: \textit{downloads} acumulados na plataforma Hugging Face, em
milhões.\end{minipage}\par\smallskip
\textit{Fonte: Elaboração própria, com dados de \citeonline{Radford2019},
\citeonline{Rajbhandari2020}, \citeonline{Frantar2023}, \citeonline{Sun2024},
\citeonline{Gao2020} e da plataforma Hugging Face (acesso em 28 set. 2026).}
\end{table}

As quatro variantes do GPT-2 atendem a C5 porque \citeonline{Radford2019} publicam, na
Tabela 3 do artigo original, a \textit{perplexity} de cada uma no WikiText-2: 29,41
(\textit{small}), 22,76 (\textit{medium}), 19,93 (\textit{large}) e 18,34 (XL). A relação
entre o valor publicado para o \textit{small} e o medido neste trabalho, obtidos por
protocolos distintos, é discutida na Seção~\ref{sec:res_protocolo}.

O GPT-2 \textit{small} acumula, além disso, 930,1 milhões de \textit{downloads} na
plataforma Hugging Face, o maior valor entre os candidatos de geração de texto considerados
--- o BERT-base tem mais \textit{downloads}, mas é um codificador, já excluído por C1. Nos
30 dias anteriores ao acesso, o modelo somou 15,4 milhões de \textit{downloads}, o segundo
maior volume entre todos os modelos de geração de texto da plataforma (acesso em 28 set.
2026). Esse dado reforça a representatividade do modelo, sem substituir os critérios de
seleção.

Dos nove candidatos, seis não atendem a todos os critérios. O BERT-base cai em C1: sua
arquitetura é apenas codificadora, pré-treinada com modelagem de linguagem mascarada
\cite{Devlin2019}, de modo que a \textit{perplexity} autorregressiva não se aplica --- ela
não é bem definida para modelos de linguagem mascarada \cite{HuggingFacePerplexity} ---,
o que o faz cair também em C5; e cai em D4, por ser pré-treinado sobre a Wikipédia em
inglês \cite{Devlin2019}. O Llama 2 7B cai em C2:
o acesso a seus pesos na plataforma Hugging Face está condicionado a aprovação manual
(verificado em 28 set. 2026). Cai também em C3 --- seus 6,74 bilhões de parâmetros exigiriam
107,8 GB para o \textit{fine-tuning} completo, muito acima dos 16 GB disponíveis. Atende,
porém, a C5: \citeonline{Sun2024} reportam \textit{perplexity} de 5,12 no WikiText-2 para o
modelo denso. O GPT-2 XL
cai isoladamente em C3: \citeonline{Rajbhandari2020} apontam que um modelo de 1,5 bilhão de
parâmetros exige, no mínimo, 24 GB de memória para o treinamento em precisão mista com
Adam, acima do limite disponível. O DistilGPT2 cai em C4, por ser obtido por destilação de
conhecimento supervisionada pelo GPT-2 de 124M \cite{Sanh2019}.

O OPT-125M e o Pythia-160M exigem exame mais detido, pois atendem a C1--C4 e C6. O
OPT-125M atende também a C5: \citeonline{Frantar2023} reportam \textit{perplexity} de 27,66
para o modelo denso no WikiText-2, em estudo de poda. Ambos, porém, caem em D4 em relação
ao WikiText-2, pois incluem a Wikipédia no pré-treinamento: o OPT foi pré-treinado sobre
uma seleção de componentes do The Pile que inclui a Wikipédia \cite{Zhang2022}, e o Pythia
sobre o The Pile completo \cite{Biderman2023}, cujo componente Wikipedia (en) soma 6,38
GiB \cite{Gao2020}. O GPT-2, ao contrário, teve os documentos da Wikipédia removidos de
sua base de dados de pré-treinamento \cite{Radford2019} (Seção~\ref{subsec:dataset_justificativa}).
Nenhum dos dois modelos forma par admissível com as demais bases de dados candidatas
(Tabela~\ref{tab:candidatos_dataset}): o WikiText-103 cai em D4 pelo mesmo motivo; o Penn
Treebank e o OpenWebText caem em critérios que independem do modelo (D2 e D3,
respectivamente); e The Pile cai em D4, por ser a base de dados de pré-treinamento do Pythia e
fornecer parte do pré-treinamento do OPT.

Restam três
candidatos que atendem a todos os critérios: as variantes \textit{small},
\textit{medium} e \textit{large} do GPT-2, esta última com margem estreita em C3 --- dos
12,4 GB estimados, restam cerca de 3,6 GB para as ativações. Aplicou-se, então, a regra de
desempate da Seção~\ref{subsec:metodologia_selecao_modelo}: entre os candidatos que atendem a todos os
critérios, adota-se o de menor número de parâmetros. O GPT-2 \textit{medium} e o
\textit{large} têm, respectivamente, cerca de 2,9 e 6,2 vezes mais parâmetros que o
\textit{small}; a diferença é relevante porque a campanha experimental deste trabalho
registrou 195 execuções completas de avaliação --- 20 da poda por magnitude, 60 da poda
estruturada, 48 do perfil de sensibilidade nas duas taxas-sonda, 60 da alocação por
sensibilidade e 7 de \textit{fine-tuning} de recuperação ---, com resultados reunidos em
\texttt{codigo/resultados/}.

No ambiente adotado, uma época de \textit{fine-tuning} do GPT-2 \textit{small} sobre o
conjunto de treino do WikiText-2 (4.743 blocos de 512 \textit{tokens}, isto é, 2.428.416
\textit{tokens}) levou 274 s, cerca de 4,6 minutos, o que corresponde a aproximadamente 8,9
mil \textit{tokens} por segundo; os sete \textit{fine-tunings} da campanha somaram cerca de
31 minutos. Supondo custo proporcional ao número de parâmetros --- aproximação coerente com
a contagem de FLOPs da Tabela~\ref{tab:gpt2_config}, em que 126,6 GFLOPs por janela de 512
\textit{tokens} correspondem a cerca de 2 operações por parâmetro por \textit{token} ---, o
mesmo conjunto de \textit{fine-tunings} levaria cerca de 1,5 hora com o \textit{medium} e
3,2 horas com o \textit{large}; trata-se de estimativas, não de medições diretas. A escolha
recaiu, portanto, sobre o GPT-2 \textit{small}.

Cabe observar que, embora seja a menor variante da família GPT-2, o modelo é cerca de
cinco vezes maior que a ResNet-50, arquitetura tradicionalmente usada em estudos de poda
em visão computacional. Essa diferença de escala é justamente o que torna o trabalho
pertinente: as conclusões da literatura de poda foram estabelecidas majoritariamente em
redes convolucionais, e não é evidente que se transfiram para modelos com esse perfil de
parâmetros.

\subsection{Seleção da Base de Dados}\label{sec:dataset}

\subsubsection{O WikiText-2}\label{subsec:wikitext2}

O WikiText-2 é uma base de dados de língua inglesa extraída de artigos verificados da
Wikipédia, proposta por \citeonline{Merity2017}, com \textit{perplexity} do GPT-2 nessa base
de dados reportada por \citeonline{Radford2019}. Diferentemente de bases de dados construídas a partir de
frases isoladas, o WikiText preserva a estrutura
original dos artigos --- parágrafos completos, pontuação e capitalização mantidas ---, o
que o torna adequado para medir a capacidade de um modelo de acompanhar contexto ao longo
de textos extensos.

Empregou-se especificamente a variante \texttt{wikitext-2-raw-v1}, na qual o texto não
sofre substituição de palavras raras por um marcador genérico. Essa escolha é necessária
porque o GPT-2 utiliza tokenização por \textit{subwords}, capaz de representar palavras
fora do vocabulário decompondo-as em fragmentos; a variante processada anularia essa
capacidade e distorceria a perplexity medida.

\subsubsection{Justificativa da Escolha}\label{subsec:dataset_justificativa}

A Seção~\ref{subsec:metodologia_base} estabeleceu quatro critérios objetivos (D1 a D4)
para a seleção da base de dados de avaliação. A Tabela~\ref{tab:candidatos_dataset} aplica esses
critérios a cinco candidatos: o WikiText-2 e o WikiText-103 \cite{Merity2017}, o Penn
Treebank \cite{Merity2017}, o OpenWebText \cite{Gokaslan2019} e The Pile \cite{Gao2020}.
Como D1 e D4 dependem do modelo, a tabela os verifica para a família GPT-2, adotada na
Seção~\ref{sec:modelo}; a verificação recíproca, para os demais candidatos a modelo,
consta na Tabela~\ref{tab:candidatos_modelo}.

\begin{table}[htb]
\centering
\caption{Candidatas a base de dados de avaliação e verificação dos critérios de seleção.}
\label{tab:candidatos_dataset}
\footnotesize
\begin{tabular}{l r c c c c}
\hline
\textbf{Base de dados} & \textbf{Treino} & \textbf{D1} & \textbf{D2} & \textbf{D3} & \textbf{D4} \\
\hline
WikiText-2    & 2,1 mi \textit{tokens}   & \checkmark & \checkmark & \checkmark & \checkmark \\
WikiText-103  & 103,2 mi \textit{tokens} & \checkmark & \checkmark & \checkmark & \checkmark \\
Penn Treebank & 0,93 mi \textit{tokens}  & \checkmark & $\times$   & \checkmark & \checkmark \\
OpenWebText   & 39,8 GB                  & ---        & \checkmark & $\times$   & $\times$   \\
The Pile      & 825 GiB                  & \checkmark & \checkmark & \checkmark & $\times$   \\
\hline
\end{tabular}
\par\smallskip
\begin{minipage}{0.92\textwidth}\centering D1: resultado de referência; D2: texto sem
pré-processamento destrutivo; D3: divisão oficial; D4: ausência de sobreposição com o
pré-treinamento (critérios definidos na Seção~\ref{subsec:metodologia_base}). D1 e D4
verificados para a família GPT-2. ---: não verificado, candidato já eliminado por outro
critério.\end{minipage}\par\smallskip
\textit{Fonte: Elaboração própria, com dados de \citeonline{Merity2017},
\citeonline{Radford2019}, \citeonline{Gao2020} e da plataforma Hugging Face (acesso em 28
set. 2026).}
\end{table}

Dos cinco candidatos, três não atendem a todos os critérios. O Penn Treebank cai em D2: na
versão usual, as palavras estão em caixa baixa, a pontuação foi removida e o vocabulário
foi limitado a 10 mil palavras \cite{Merity2017}, ao passo que o WikiText mantém pontuação,
caixa original e números \cite{Merity2017}, e a variante \textit{raw} adotada neste
trabalho não substitui palavras raras. O OpenWebText cai em D3, pois é distribuído com uma
única partição, de treino, sem divisão oficial em validação e teste (8.013.769 documentos,
39,8 GB, verificado na plataforma Hugging Face); e em D4, por ser uma reprodução aberta do
WebText, a base de dados de pré-treinamento do GPT-2 \cite{Gokaslan2019}. The Pile atende a D1 ---
\citeonline{Gao2020} avaliam as quatro variantes do GPT-2 em seu conjunto de teste, com a
\textit{perplexity} convertida em bits por \textit{byte} --- e a D3, por dispor de
partições oficiais de treino, validação e teste --- validação e teste com cerca de 0,1\% dos
dados cada, mais de 1 GiB cada \cite{Gao2020}; cai em D4 porque inclui o OpenWebText2 (62,77
GiB), construído pelo mesmo método do WebText, com votos no Reddit como filtro de qualidade
dos links \cite{Gao2020}. Em sentido oposto, os autores do GPT-2 removeram do WebText todos
os documentos da Wikipédia justamente para evitar sobreposição com conjuntos de avaliação
\cite{Radford2019}; a sobreposição de 8-gramas medida por eles entre os conjuntos de teste e
o treino do WebText é de 1,63\% no WikiText-2, 2,42\% no WikiText-103 e 0,88\% no Penn
Treebank (Tabela 6 de \citeonline{Radford2019}).

Restam dois candidatos que atendem a todos os critérios: o WikiText-2 e o WikiText-103.
Ambos compartilham os mesmos conjuntos de validação e teste, com 217.646 e 245.569
\textit{tokens}, respectivamente \cite{Merity2017}; diferem apenas no conjunto de treino,
cerca de 49 vezes maior no WikiText-103. Como o conjunto de teste é o mesmo, a avaliação de
\textit{perplexity} é idêntica nos dois; a diferença está inteiramente no custo do
\textit{fine-tuning}, determinado pelo treino. Com a vazão medida na
Seção~\ref{subsec:justificativa_gpt2} --- uma época sobre o treino do WikiText-2 em cerca de
4,6 minutos ---, uma época sobre o WikiText-103 levaria cerca de 3,8 horas por
\textit{fine-tuning}, e os sete \textit{fine-tunings} da campanha passariam de cerca de 31
minutos para cerca de 25 horas (estimativa proporcional ao número de \textit{tokens} de
treino). Aplicou-se, então, a regra de desempate da Seção~\ref{subsec:metodologia_base}:
adotou-se o de menor conjunto de treino, o WikiText-2.

Registre-se que a base de dados é parametrizável na implementação: a substituição pelo
WikiText-103, de mesma natureza e cerca de cinquenta vezes maior, exige apenas a alteração
de um argumento de linha de comando. Essa flexibilidade foi mantida deliberadamente para
que trabalhos futuros possam verificar a generalidade dos resultados sem reescrever o
código.

\subsubsection{Pré-processamento e Tokenização}\label{subsec:preprocessamento}

O pré-processamento foi intencionalmente mínimo, para que a perplexity medida refletisse a
base de dados e não decisões de limpeza. As linhas do conjunto foram unidas por quebra dupla,
preservando a separação entre parágrafos tal como na base de dados original, e o texto resultante
foi tokenizado como um fluxo contínuo único.

A tokenização empregou o tokenizador original do GPT-2, baseado no algoritmo BPE
(\textit{Byte Pair Encoding}), com vocabulário de 50.257 \textit{tokens}. Manter o
tokenizador original é indispensável: os \textit{embeddings} do modelo foram aprendidos
para esse vocabulário específico, e qualquer substituição invalidaria os pesos
pré-treinados.

\subsubsection{Divisão dos Dados}\label{subsec:divisao_dados}

Adotou-se a divisão oficial da base de dados. O \textbf{conjunto de teste} foi utilizado em todas as
medições de perplexity reportadas, tanto do modelo base quanto de todas as configurações
comprimidas. O \textbf{conjunto de treinamento} foi utilizado exclusivamente na etapa de
\textit{fine-tuning} de recuperação (Seção~\ref{sec:finetuning}). O \textbf{conjunto de
validação} foi empregado apenas para acompanhar a evolução do treinamento ao longo dos
passos, quando essa curva foi solicitada, e nunca para selecionar resultados reportados.

Essa separação estrita importa: como as varreduras \textit{one-shot} não realizaram
treinamento algum, elas não tiveram contato com o conjunto de treinamento, e a perplexity de
teste é uma medida limpa do dano causado pela poda.

\subsection{Métricas de Avaliação}\label{sec:metricas}

Toda configuração --- o modelo base, as podadas pelos métodos da literatura e as podadas
pela estratégia otimizada --- foi avaliada pelo mesmo conjunto de métricas, cobrindo as três
dimensões definidas na Seção~\ref{sec:metodologia_metricas}. A
Tabela~\ref{tab:metricas} as resume.

\begin{table}[htb]
\centering
\caption{Métricas de avaliação coletadas para cada configuração.}
\label{tab:metricas}
\small
\begin{tabular}{p{3.2cm} p{6.1cm} p{1.7cm} p{2.2cm}}
\hline
\textbf{Métrica} & \textbf{O que mede} & \textbf{Unidade} & \textbf{Ferramenta} \\
\hline
Perplexity & Qualidade da modelagem da linguagem; quanto menor, melhor & --- & PyTorch \\
Parâmetros totais & Parâmetros existentes na arquitetura & unidades & PyTorch \\
Parâmetros não nulos & Parâmetros efetivamente diferentes de zero & unidades & PyTorch \\
Esparsidade real & Fração de parâmetros podáveis eliminados & \% & PyTorch \\
Tamanho do modelo & Memória ocupada pelos parâmetros & MB & PyTorch \\
FLOPs & Operações de ponto flutuante por inferência & FLOPs & \texttt{thop} \\
Tempo de inferência & Tempo médio de uma passagem direta & ms & \texttt{perf\_counter} \\
Energia consumida & Energia gasta durante a avaliação & kWh & \texttt{CodeCarbon} \\
Emissões de $CO_2$ & Emissões equivalentes ao consumo & kg & \texttt{CodeCarbon} \\
\hline
\end{tabular}\\
\textit{Fonte: Elaboração própria.}
\end{table}

\subsubsection{Perplexity}\label{subsec:metrica_perplexity}

A perplexity é a métrica de desempenho preditivo adotada. Ela foi calculada pelo método da
janela deslizante, padrão para modelos autorregressivos: o fluxo de \textit{tokens} do
conjunto de teste foi percorrido em janelas do tamanho máximo de contexto do modelo, com
passo menor que a janela, de modo que cada \textit{token} fosse previsto dispondo de
contexto anterior substancial. Apenas os \textit{tokens} ainda não avaliados contribuíram
para a perda em cada janela; os demais entraram como contexto.

% AJUSTE: 21/10 — o "compromisso usual na literatura" ganha referência
% (HuggingFacePerplexity usa a mesma janela de 1.024 com passo de 512 no GPT-2).
Empregou-se janela de 1.024 \textit{tokens} com passo de 512. Passos menores produzem
estimativas mais precisas ao custo de mais computação; o valor adotado é o compromisso
usual na literatura \cite{HuggingFacePerplexity}, e foi mantido idêntico em todas as
configurações avaliadas.

\subsubsection{Custo Computacional}\label{subsec:metrica_custo}

% AJUSTE: 21/10 — "elas coincidem" na poda estruturada podia ser lido como comparação
% com o modelo de referência; explicita-se que coincidem as duas contagens do modelo
% podado, ambas abaixo das do modelo de referência.
A contagem de \textbf{parâmetros} é reportada em duas formas complementares. O total
reflete o que existe na arquitetura; a contagem de não nulos reflete o que sobrou depois da
poda. Na poda não estruturada as duas divergem --- há zeros, mas a arquitetura é a mesma.
Na poda estruturada, as duas contagens do modelo podado coincidem, e ambas ficam abaixo
das do modelo base: as estruturas removidas deixam de existir, em vez de
permanecerem como zeros. Essa divergência é, por si só, um resultado: ela evidencia a
diferença entre compressão contábil e compressão física.

A \textbf{esparsidade real} reportada nas tabelas foi sempre calculada sobre o conjunto
podável, mas a natureza do cálculo acompanha a estratégia: na poda por magnitude é a fração
de pesos anulados nas 48 matrizes; na poda estruturada é a fração de parâmetros que
deixaram de existir nas mesmas projeções, vieses incluídos, já que a remoção de uma
estrutura elimina também o viés que lhe corresponde. As duas grandezas medem a mesma coisa
--- quanto do conjunto podável foi eliminado.

% AJUSTE: 21/10 — o parágrafo dos FLOPs era vago ("camadas lineares não convencionais")
% e afirmava que a omissão da atenção era idêntica entre estratégias, o que não vale para
% a poda estruturada. Reescrito sem nomes de biblioteca, com a decomposição medida:
% projeções 86,97 GFLOPs, cabeça de linguagem 39,52, LayerNorm 0,08 (total 126,58, igual
% ao CSV); atenção não contada 9,66 GFLOPs a 512 tokens (7,1% de 136,24). "Contagem pelo
% grafo" passa a "contagem por camadas", em coerência com a nova descrição.
% AJUSTE: 24/10 — "camada" passou a designar os 12 blocos (Seção do GPT-2); aqui o termo
% designava as unidades de software percorridas pela ferramenta. Passa a "módulo".
Os \textbf{FLOPs} foram contabilizados por uma ferramenta que percorre os módulos do
modelo --- as unidades que o compõem na implementação, como as projeções lineares --- e
soma o custo de cada um, segundo uma regra associada ao tipo do módulo. Em um módulo
linear, o custo é de duas operações, uma multiplicação e uma soma, por peso da matriz,
para cada \textit{token} processado. No GPT-2, as projeções das camadas efetuam essa mesma
operação, mas são implementadas por um tipo de módulo que a ferramenta não reconhece e ao
qual atribui custo nulo. Sem tratamento, a contagem ignoraria as matrizes que a poda
modifica, responsáveis por 87,0 dos 126,6 GFLOPs do modelo base, e permaneceria
constante sob qualquer poda. Por isso, aplicou-se a essas projeções a mesma regra de
contagem dos módulos lineares.

As multiplicações de matriz internas ao mecanismo de atenção não são executadas por um
módulo próprio, e sim diretamente no cálculo da atenção, e não foram contabilizadas em
nenhuma configuração. Elas representam cerca de 7\% do custo do
modelo base. Na poda por magnitude, essa parcela é constante; na poda
estruturada, ela diminui com a remoção de cabeças de atenção, de modo que a redução de
FLOPs reportada é ligeiramente inferior à real.

A implementação prevê ainda um estimador analítico de
reserva --- duas operações por parâmetro não pertencente aos \textit{embeddings}, por
\textit{token} processado ---, acionado apenas se a contagem por módulos falhar. Todos os
valores reportados no Capítulo~\ref{chap:resultados} provêm da contagem por módulos, o que
se verifica pelo fato de os FLOPs variarem com a poda estruturada; o estimador de reserva
produziria uma grandeza de natureza diferente e não deve ser misturado a eles.

O \textbf{tempo de inferência} foi medido sobre uma entrada representativa de 512
\textit{tokens}, com execuções de aquecimento descartadas --- para mitigar efeitos de
inicialização da GPU --- seguidas de execuções cronometradas, das quais se reporta média e
desvio-padrão.

\subsubsection{Impacto Ambiental}\label{subsec:metrica_energia}

O consumo de \textbf{energia} foi medido com o \texttt{CodeCarbon}, que monitora o
\textit{hardware} durante a execução do bloco de avaliação e reporta a energia consumida em
quilowatts-hora. As \textbf{emissões de $CO_2$} foram estimadas a partir desse consumo e da
intensidade de carbono da rede elétrica onde o código executa.

Sobre esta última métrica cabe uma ressalva importante, identificada durante os
experimentos: a intensidade de carbono utilizada variou conforme a sessão de execução. Em
consequência, as emissões não são comparáveis entre execuções realizadas em sessões
distintas, e a análise dos resultados apoia-se no consumo de energia, não nas emissões
derivadas dele.

\section{Etapa 2 --- Poda sem Otimização}\label{sec:metodos_poda}

Na segunda etapa, aplicaram-se duas estratégias clássicas de poda: a poda não estruturada
por magnitude \cite{Han2015} e a poda estruturada por norma $L_1$ \cite{Li2017}, esta
última adaptada ao \textit{Transformer}. Cada estratégia foi aplicada de forma
independente e sempre a partir do modelo base, em toda a faixa de
esparsidade --- de 10\% a 90\%, em passos de 10 pontos percentuais --- e sob dois escopos
de alocação distintos. A escolha dessas duas estratégias é deliberada: elas ocupam os dois
extremos do espectro de granularidade e, como se verá, comportam-se de forma oposta quanto
ao impacto real sobre o \textit{hardware}.

\subsection{Seleção das Estratégias}\label{subsec:selecao_estrategias}

A Seção~\ref{sec:metodologia_poda_sem} estabeleceu três critérios (P1 a P3) para a escolha
das estratégias de poda sem otimização. A Tabela~\ref{tab:candidatos_poda} aplica esses
critérios à poda por magnitude \cite{Han2015}, à poda estruturada por norma $L_1$
\cite{Li2017} e a quatro alternativas da literatura.

\begin{table}[htb]
\centering
\caption{Estratégias de poda candidatas e verificação dos critérios de seleção.}
\label{tab:candidatos_poda}
\footnotesize
\setlength{\tabcolsep}{3.5pt}
\begin{tabular}{l l l c r c c}
\hline
\textbf{Estratégia} & \textbf{Granularidade} & \textbf{Critério de importância} & \textbf{Ano} & \textbf{Citações} & \textbf{P1} & \textbf{P2} \\
\hline
Magnitude                  & fina   & $|w|$ de cada peso                          & 2015 & 8.047 & \checkmark & \checkmark \\
Norma $L_1$                & grossa & norma $L_1$ dos pesos da unidade            & 2017 & 4.214 & \checkmark & \checkmark \\
SparseGPT                  & fina   & reconstrução com dados de calibração        & 2023 & 1.552 & $\times$   & \checkmark \\
Wanda                      & fina   & $|w|$ $\times$ norma das ativações          & 2024 &   999 & $\times$   & \checkmark \\
\textit{Movement pruning}  & fina   & movimento dos pesos no \textit{fine-tuning} & 2020 &   654 & $\times$   & $\times$   \\
Cabeças por gradiente      & grossa & gradiente da perda em relação à máscara     & 2019 & 1.513 & $\times$   & \checkmark \\
\hline
\end{tabular}
\par\smallskip
\begin{minipage}{0.92\textwidth}\centering P1: importância calculada apenas a partir dos
pesos; P2: aplicação \textit{one-shot} (critérios definidos na
Seção~\ref{sec:metodologia_poda_sem}). P3 refere-se ao conjunto de estratégias e é
discutido no texto. Ano: ano de publicação do artigo que propôs a estratégia. Citações:
citações desse artigo na base Semantic Scholar, indicador complementar de
representatividade, que não integra os critérios.\end{minipage}\par\smallskip
\textit{Fonte: Elaboração própria, com base em \citeonline{Han2015}, \citeonline{Li2017},
\citeonline{Frantar2023}, \citeonline{Sun2024}, \citeonline{Sanh2020} e
\citeonline{Michel2019}, e com dados da base Semantic Scholar (acesso em 3 out. 2026).}
\end{table}

As quatro alternativas caem em P1, pois todas estimam a importância com base em dados. O
SparseGPT e o Wanda, referências recentes da poda não estruturada de LLMs, utilizam 128
sequências de calibração extraídas da base de dados C4: o primeiro para reconstruir, com
informação de segunda ordem, os pesos remanescentes de cada camada \cite{Frantar2023}; o
segundo para ponderar a magnitude de cada peso pela norma da ativação de entrada
correspondente \cite{Sun2024}. A importância das cabeças em \citeonline{Michel2019} é
estimada pelo gradiente da perda sobre uma base de dados. A \textit{movement pruning}
cai também em P2, pois a importância de cada peso é aprendida ao longo do
\textit{fine-tuning} \cite{Sanh2020}. A exclusão dessas estratégias não decorre de
desempenho inferior: elas já incorporam o tipo de informação que a análise de
sensibilidade introduz, e sua combinação com a alocação por sensibilidade fica fora do
escopo deste trabalho. Restam a poda por magnitude e a poda estruturada por norma $L_1$,
que atendem a P1 e P2 e, juntas, atendem a P3: a primeira opera sobre pesos individuais e
a segunda sobre unidades arquiteturais inteiras.

% AJUSTE: revisão 30/09 — representatividade medida por citações, a pedido do orientador.
Quanto ao indicador complementar de representatividade, os artigos que propuseram as duas
estratégias adotadas são os mais citados entre as seis candidatas: 8.047 citações para
\citeonline{Han2015} e 4.214 para \citeonline{Li2017}, contra 1.552 do SparseGPT, o
terceiro (Semantic Scholar, acesso em 3 out. 2026). Como as citações se acumulam com o
tempo, a diferença reflete em parte a antiguidade dos trabalhos --- as quatro alternativas
foram publicadas entre 2019 e 2024. O dado reforça a condição de referência das duas
estratégias na literatura de poda, sem substituir os critérios de seleção.
% AJUSTE: 05/10 — frase trazida da Metodologia (Etapa 2), que deixou de nomear as
% estratégias escolhidas.
A mesma condição é corroborada pela meta-análise de \citeonline{Blalock2020}, que reuniu os
resultados de 81 trabalhos de poda e aponta como um dos resultados consistentes da área que
a poda pela magnitude dos pesos comprime substancialmente as redes sem reduzir a acurácia.

\subsection{Parâmetros Podáveis}\label{subsec:parametros_podaveis}

% AJUSTE: 21/10 — (i) a "convenção consolidada na literatura" ganha referência: SparseGPT
% (Frantar2023, Seç. 4: "prune all linear layers, excluding the embeddings and the head, as
% standard") e Wanda (Sun2024, Seç. 4: "skipping the first embedding layer and the final
% classification head"); (ii) "parcela expressiva" dos embeddings vira número: wte + wpe =
% 39.383.808 de 124.439.808 parâmetros (31,6%).

Uma decisão de implementação que condiciona todos os resultados é a definição de quais
parâmetros são elegíveis à poda. Adotaram-se, por convenção consolidada na literatura
\cite{Frantar2023,Sun2024}, apenas as matrizes de peso das projeções lineares internas às
camadas, quatro por camada:

\begin{itemize}
    \item \texttt{attn.c\_attn} --- projeção que produz consulta, chave e valor;
    \item \texttt{attn.c\_proj} --- projeção de saída da atenção;
    \item \texttt{mlp.c\_fc} --- primeira transformação linear do MLP;
    \item \texttt{mlp.c\_proj} --- segunda transformação linear do MLP.
\end{itemize}

% AJUSTE: 24/10 — os três itens excluídos diziam o que cada componente é, mas não por que
% fica fora da poda. Cada item ganha a razão: (i) embeddings são consulta a tabela, sem
% multiplicação de matrizes, e não contêm cabeças nem neurônios; (ii) lm_head é o mesmo
% tensor de wte (weight tying), custa 39,52 GFLOPs (decomposição de 21/10) e podá-la
% espalharia o efeito pela entrada e pela saída, o que quebra a atribuição por camada;
% (iii) vieses + LayerNorm = 82.944 + 38.400 = 121.344 parâmetros (0,098% de 124.439.808),
% LayerNorm 0,08 GFLOPs; sensibilidade ancorada em Kovaleva2021 (Findings ACL 2021, que
% inclui GPT-2). "Numericamente baratos e desproporcionalmente sensíveis" sai por ser vago.
Com 12 camadas, são 48 matrizes podáveis (Figura~\ref{fig:parametros_podaveis}). Ficaram
deliberadamente fora do escopo da poda os componentes a seguir, cada um por uma razão
própria:

\begin{itemize}
    \item as matrizes de \textit{embedding} de \textit{token} (\texttt{wte}) e de posição
          (\texttt{wpe}). Seu uso não envolve multiplicação de matrizes: a cada
          \textit{token}, o modelo apenas copia as linhas correspondentes ao \textit{token}
          e à sua posição. Anular entradas dessas matrizes reduziria o armazenamento, mas
          não o custo computacional, cuja redução é o objetivo deste trabalho. Além disso,
          elas não contêm cabeças de atenção nem neurônios do MLP, as unidades da poda
          estruturada: remover uma linha eliminaria um \textit{token} do vocabulário ou uma
          posição da sequência, e remover uma coluna alteraria a dimensão das
          representações (768), que é a dimensão de entrada ou de saída de todas as 48
          matrizes podáveis. Restringir as duas estratégias ao mesmo conjunto de matrizes
          mantém comparáveis as suas taxas de poda;
    \item a cabeça de linguagem (\texttt{lm\_head}). Diferentemente dos
          \textit{embeddings}, ela efetua computação --- 39,5 dos 126,6 GFLOPs do modelo
          base ---, mas, no GPT-2, ela e \texttt{wte} são o mesmo tensor de pesos,
          usado na entrada para representar os \textit{tokens} e na saída para pontuá-los.
          Um peso anulado na cabeça seria anulado também no \textit{embedding}, e a poda
          alteraria simultaneamente a entrada e a saída do modelo, o que impediria atribuir
          seu efeito a uma única camada, premissa da análise de sensibilidade por camada
          (Seção~\ref{subsec:fase1}). A cabeça tampouco contém as unidades da poda
          estruturada: cada uma de suas linhas corresponde a um \textit{token} do
          vocabulário. Seu custo permanece, portanto, constante em todas as configurações;
    \item vieses e parâmetros de normalização de camada. Somados, são 121.344 parâmetros,
          menos de 0,1\% do modelo, e as operações que efetuam são somas e produtos
          elemento a elemento, de custo desprezível diante das multiplicações de matriz: a
          normalização responde por 0,08 dos 126,6 GFLOPs do modelo base. Mesmo
          eliminados por inteiro, não reduziriam o custo de forma mensurável. O risco, em
          contrapartida, é alto: um parâmetro de normalização multiplica uma dimensão
          inteira da entrada da atenção ou do MLP, enquanto um peso de matriz é apenas uma
          entre 768 ou 3.072 parcelas de uma soma. \citeonline{Kovaleva2021} mostraram que
          anular os parâmetros de normalização associados a uma única dimensão --- menos de
          0,0001\% dos pesos do modelo --- degrada significativamente o desempenho de
          modelos pré-treinados, efeito também identificado no GPT-2. Na poda estruturada, o
          viés de uma estrutura removida sai junto com ela
          (Seção~\ref{subsec:metrica_custo}), mas os vieses não participam da seleção.
\end{itemize}

% AJUSTE: 21/10 — figura nova: a leitura "4 blocos" / "12 blocos dentro de 12 blocos" é
% possível sem ela. Mostra que os 12 blocos são todos podados, que as 4 matrizes estão
% dentro de cada bloco e que cabeças e neurônios são fatias dessas matrizes.
% AJUSTE: 24/10 — os rótulos "(12 cabeças)" em c_attn e "(3.072 neurônios)" em c_fc
% sugeriam que cada estrutura mora em uma só matriz. A legenda passa a dizer que cada uma
% atravessa duas projeções (cabeça: 192 colunas de c_attn + 64 linhas de attn.c_proj;
% neurônio: coluna de c_fc + linha de mlp.c_proj, como em prune_structured.py) e dá os
% totais do modelo: 12 x 12 = 144 cabeças, 12 x 3.072 = 36.864 neurônios.
\begin{figure}[htb]
\centering
\caption{Localização dos parâmetros podáveis no GPT-2 \textit{small}. O modelo empilha 12
blocos \textit{Transformer} de estrutura idêntica, aqui denominados camadas e numerados de
0 a 11; a estrutura de cada camada está detalhada à direita. Em cada camada, as quatro
matrizes de projeção, em destaque, são podáveis, o que totaliza 48 matrizes. As cabeças de
atenção e os neurônios do MLP não são matrizes à parte, e sim fatias que atravessam duas
projeções. Cada uma das 12 cabeças de uma camada ocupa 192 colunas de \texttt{attn.c\_attn} (64 para
consulta, 64 para chave e 64 para valor) e as 64 linhas de \texttt{attn.c\_proj} que
recebem sua saída. Cada um dos 3.072 neurônios ocupa uma coluna de \texttt{mlp.c\_fc}, que
o produz, e uma linha de \texttt{mlp.c\_proj}, que o consome. A poda estruturada remove a
fatia inteira, nas duas matrizes. No modelo inteiro, são 144 cabeças e 36.864 neurônios.
Os componentes tracejados ficam fora da poda.}
\label{fig:parametros_podaveis}
\begin{tikzpicture}[font=\scriptsize,
  podavel/.style={rectangle,rounded corners=2pt,draw,thick,fill=gray!35,
                  minimum width=6.4cm,minimum height=0.8cm,text width=6.1cm,align=center},
  fora/.style={rectangle,rounded corners=2pt,draw,dashed,fill=white,
               minimum width=4.2cm,minimum height=0.6cm,text width=4.0cm,align=center},
  foralargo/.style={fora,minimum width=6.4cm,text width=6.1cm},
  bloco/.style={rectangle,rounded corners=2pt,draw,thick,fill=gray!12,
                minimum width=4.2cm,minimum height=0.6cm,text width=4.0cm,align=center},
  seta/.style={-{Stealth}}]

  % --- modelo completo (esquerda) ---
  \node (in) at (0,0) {\textit{tokens} de entrada};
  \node[fora] (emb) at (0,-0.95) {\textit{Embeddings} de \textit{token}\\ e de posição};
  \node[bloco] (b1) at (0,-2.1) {\textbf{Camada 0}};
  \node[bloco] (b2) at (0,-3.1) {\textbf{Camada 1}};
  \node (dots) at (0,-3.85) {$\vdots$};
  \node[bloco] (b12) at (0,-4.6) {\textbf{Camada 11}};
  \node[fora] (lnf) at (0,-5.6) {Normalização final};
  \node[fora] (head) at (0,-6.75) {Cabeça de linguagem\\ (pesos do \textit{embedding})};
  \node (out) at (0,-7.8) {próximo \textit{token}};
  \draw[seta] (in) -- (emb);
  \draw[seta] (emb) -- (b1);
  \draw[seta] (b1) -- (b2);
  \draw (b2) -- (dots);
  \draw[seta] (dots) -- (b12);
  \draw[seta] (b12) -- (lnf);
  \draw[seta] (lnf) -- (head);
  \draw[seta] (head) -- (out);

  % --- estrutura de um bloco (direita) ---
  \node[foralargo] (ln1) at (7.6,-0.6) {Normalização de camada};
  \node[podavel] (cattn) at (7.6,-1.75) {\texttt{attn.c\_attn} \quad 768 $\rightarrow$ 2.304\\
        consulta, chave e valor (12 cabeças)};
  \node[podavel] (cproj) at (7.6,-2.95) {\texttt{attn.c\_proj} \quad 768 $\rightarrow$ 768\\
        saída da atenção};
  \node[foralargo] (ln2) at (7.6,-4.1) {Normalização de camada};
  \node[podavel] (cfc) at (7.6,-5.25) {\texttt{mlp.c\_fc} \quad 768 $\rightarrow$ 3.072\\
        expansão (3.072 neurônios)};
  \node[podavel] (mproj) at (7.6,-6.45) {\texttt{mlp.c\_proj} \quad 3.072 $\rightarrow$ 768\\
        retorno à dimensão 768};
  \draw[seta] (ln1) -- (cattn);
  \draw[seta] (cattn) -- (cproj);
  \draw[seta] (cproj) -- (ln2);
  \draw[seta] (ln2) -- (cfc);
  \draw[seta] (cfc) -- (mproj);
  \node[draw,dashed,gray!70,rounded corners=3pt,inner sep=8pt,fit=(ln1)(mproj),
        label={[font=\scriptsize,text=black]above:\textbf{Cada camada} (bloco \textit{Transformer}; estrutura idêntica nas 12)}]
        (frame) {};
  \draw[decorate,decoration={brace,amplitude=4pt}]
        ($(cattn.north east)+(0.5,0)$) -- ($(cproj.south east)+(0.5,0)$)
        node[midway,right=5pt] {Atenção};
  \draw[decorate,decoration={brace,amplitude=4pt}]
        ($(cfc.north east)+(0.5,0)$) -- ($(mproj.south east)+(0.5,0)$)
        node[midway,right=5pt] {MLP};

  % --- ampliação da camada 0 ---
  \draw[dashed,gray!70] (b1.north east) -- (frame.north west);
  \draw[dashed,gray!70] (b1.south east) -- (frame.south west);

  % --- legenda ---
  \node[rectangle,rounded corners=1pt,draw,thick,fill=gray!35,
        minimum width=0.5cm,minimum height=0.3cm] (lg1) at (4.7,-7.75) {};
  \node[right=3pt of lg1] {podável: 4 matrizes por camada, 48 no total};
  \node[rectangle,rounded corners=1pt,draw,dashed,
        minimum width=0.5cm,minimum height=0.3cm] (lg2) at (4.7,-8.25) {};
  \node[right=3pt of lg2] {fora da poda};
\end{tikzpicture}
\end{figure}

Essa restrição tem consequência direta na leitura dos resultados: quando se afirma que uma
configuração tem 50\% de esparsidade, trata-se de 50\% das matrizes podáveis, não do
modelo inteiro. Como os \textit{embeddings} do GPT-2 concentram 31,6\% dos
parâmetros, a esparsidade sobre o total é sempre menor que a esparsidade sobre o conjunto
podável, e as duas grandezas são reportadas separadamente.

\subsection{Método 1 --- Poda Não Estruturada por Magnitude}\label{subsec:metodo_magnitude}

O primeiro método segue \citeonline{Han2015}. O critério de importância é o valor absoluto
do peso: assume-se que pesos próximos de zero contribuem pouco para a saída e podem ser
anulados. Dada uma taxa de esparsidade alvo, calculou-se o limiar correspondente ao
quantil apropriado da distribuição de $|w|$ e zeraram-se todos os pesos abaixo dele.

A implementação preserva a característica essencial do método: os pesos são
\textbf{zerados}, não removidos. As matrizes conservam suas dimensões originais e a
arquitetura permanece intacta --- o que, como discutido na Seção~\ref{sec:pruning}, é
precisamente a razão pela qual esse método não reduz o custo de inferência em
\textit{hardware} convencional.

\subsection{Método 2 --- Poda Estruturada por Norma $L_1$}\label{subsec:metodo_estruturado}

O segundo método adapta \citeonline{Li2017} ao \textit{Transformer}. O trabalho original
poda filtros convolucionais ordenados por norma $L_1$; aqui, as unidades estruturais
equivalentes são as \textbf{cabeças de atenção} e os \textbf{neurônios do MLP}, os 3.072
de sua dimensão intermediária. Para cada uma dessas estruturas calculou-se a norma $L_1$ dos pesos
que lhe correspondem, e removeram-se as de menor norma.

A diferença decisiva em relação ao método anterior é que a remoção é \textbf{física}. As
matrizes foram efetivamente reconstruídas com dimensões menores: as colunas de consulta,
chave e valor associadas às cabeças removidas foram eliminadas de \texttt{c\_attn}, as
linhas correspondentes saíram de \texttt{c\_proj}, e a contabilidade interna do módulo de
atenção --- número de cabeças e dimensão de divisão --- foi atualizada de acordo. O
resultado não é um modelo esparso, mas uma arquitetura menor e densa.

% AJUSTE: 24/10 — (i) o parágrafo não dizia qual alvo os resultados usam; passa a dizer
% (ambos, gerar_tabelas.py). (ii) "romperia o fluxo residual" era incorreto: sem cabeças,
% a atenção devolve só o viés de c_proj e o bloco vira h = x + b, com o fluxo residual
% intacto. A salvaguarda passa a ser justificada pela granularidade: bloco sem cabeças ou
% sem neurônios = subcamada removida inteira = poda de camadas, fora do escopo.
% AJUSTE: 24/10 (2) — os alvos isolados aparecem, sim, nos Resultados (comparação cabeças ×
% neurônios, Seção res_estruturada); a frase passa a dizer isso. "Bloco" → "camada", e
% "poda de camadas" sai, pois com camada = bloco leria-se como remoção de blocos inteiros.
A implementação admite três alvos: apenas cabeças, apenas neurônios do MLP, ou ambos
simultaneamente. Os resultados do Capítulo~\ref{chap:resultados} empregam os dois alvos
simultâneos, com a mesma taxa aplicada a cada tipo de estrutura; os alvos isolados
aparecem apenas na comparação entre cabeças e neurônios da
Seção~\ref{sec:res_estruturada}. Uma salvaguarda impede que qualquer camada fique sem
nenhuma cabeça ou sem nenhum neurônio: sem ela, a poda deixaria de remover estruturas e
passaria a remover subcamadas inteiras --- a atenção ou o MLP de uma camada ---,
granularidade mais grossa que a adotada e fora do escopo deste trabalho.

% AJUSTE: 24/10 (3) — pedido do Enzo: a primeira razão ganha o porquê. Pesos por estrutura
% (sem vieses): cabeça = 192 x 768 (c_attn) + 64 x 768 (attn.c_proj) = 196.608; neurônio =
% 768 (c_fc) + 768 (mlp.c_proj) = 1.536; razão 128. Por camada: 12 cabeças = 2.359.296 de
% 7.077.888 pesos podáveis (1/3); 3.072 neurônios = 4.718.592 (2/3). Conferido em
% estruturada.csv: alvo só cabeças a 50% -> esparsidade real 16,67%; só neurônios -> 33,32%;
% ambos -> 49,99%. Com os dois alvos à mesma taxa as parcelas somam a taxa, e a divergência
% vem do arredondamento das cabeças (uniforme 10%: 1 de 12 cabeças, real 9,44%).
Convém registrar que a esparsidade alvo, neste método, é expressa como fração de
\textbf{estruturas} removidas, não de parâmetros. Duas razões fazem a fração de parâmetros
efetivamente eliminada divergir da taxa solicitada, e por isso ela é reportada
separadamente em todas as tabelas. A primeira é que cabeças e neurônios têm custos
diferentes em parâmetros, já que cada estrutura é removida com todos os pesos que lhe
pertencem, e uma cabeça de atenção reúne muito mais pesos que um neurônio do MLP
(Figura~\ref{fig:parametros_podaveis}). Cada cabeça ocupa 192 colunas de
\texttt{attn.c\_attn} e 64 linhas de \texttt{attn.c\_proj}, todas com 768 pesos, o que
soma 196.608 pesos; cada neurônio ocupa apenas uma coluna de \texttt{mlp.c\_fc} e uma linha
de \texttt{mlp.c\_proj}, também com 768 pesos cada, o que soma 1.536. Remover uma única
cabeça elimina, portanto, tantos pesos quanto remover 128 neurônios. Como a taxa conta
estruturas, e não pesos, a fração de parâmetros eliminada depende de quais estruturas são
removidas: as 12 cabeças de uma camada concentram um terço de seus pesos podáveis
(2.359.296 de 7.077.888), e os 3.072 neurônios, os outros dois terços. Uma taxa de 50\%
aplicada apenas às cabeças elimina, assim, cerca de 16,7\% dos parâmetros podáveis, um
terço da taxa; aplicada apenas aos neurônios, cerca de 33,3\%, dois terços dela. Quando os
dois alvos recebem a mesma taxa, como nos resultados do Capítulo~\ref{chap:resultados}, as
duas parcelas somam a própria taxa, e a divergência passa a vir da segunda razão, descrita
a seguir, cujo efeito a diferença de tamanho amplia: cada cabeça que o arredondamento
deixa de remover conserva tantos pesos quanto 128 neurônios.

A segunda razão é a granularidade discreta das estruturas: o número de
unidades a remover é a parte inteira do produto entre a taxa e a contagem disponível, o que
nos neurônios do MLP --- 3.072 por camada --- é praticamente imperceptível, mas nas cabeças
de atenção, que são apenas 12 por camada, produz arredondamentos visíveis. Uma taxa de
90\%, por exemplo, remove 10 das 12 cabeças, ou seja, 83\%.

\subsection{Escopos de Alocação}\label{subsec:escopos}

Ambos os métodos foram avaliados sob dois escopos, que correspondem às alocações global e
uniforme definidas na Seção~\ref{sec:metodologia_poda_sem} e representam formas distintas de
distribuir a poda entre as camadas:

\begin{itemize}
    \item \textbf{Global:} o critério de importância é aplicado ao modelo inteiro de uma
          só vez. Um único limiar --- ou uma única ordenação de estruturas --- decide o
          que sai, de modo que camadas com pesos sistematicamente menores cedem
          proporcionalmente mais. Na poda estruturada com os dois alvos simultâneos, o
          escopo global compreende duas ordenações independentes, uma de cabeças e outra
          de neurônios, cada uma abrangendo as doze camadas: as normas $L_1$ de uma
          cabeça de atenção e de um neurônio do MLP não são grandezas comparáveis, pois
          agregam números distintos de pesos, e ordená-las em conjunto faria a escolha
          depender do tamanho da estrutura em vez de sua importância relativa. A norma
          $L_1$ de uma estrutura é a soma dos valores absolutos de seus pesos, e uma cabeça
          soma 196.608 pesos, 128 vezes os 1.536 de um neurônio
          (Seção~\ref{subsec:metodo_estruturado}). Ainda que os pesos das duas estruturas
          tivessem o mesmo valor típico, a norma de uma cabeça seria cerca de 128 vezes a
          de um neurônio, apenas por somar mais parcelas. É o que ocorre no GPT-2: o valor
          absoluto médio de um peso é praticamente o mesmo nas duas estruturas, cerca de
          0,1, e a cabeça de menor norma tem norma 14 vezes maior que a do neurônio de
          maior norma. Em uma ordenação conjunta, portanto, todos os neurônios viriam antes
          da primeira cabeça, e nenhuma cabeça seria removida em nenhuma das taxas
          avaliadas, até 90\%: a poda dos dois alvos se reduziria à poda apenas de
          neurônios. As duas ordenações independentes evitam esse efeito, já que cada
          estrutura passa a ser comparada somente com estruturas do mesmo tipo e do mesmo
          tamanho, e garantem que a taxa de poda incida igualmente sobre os dois tipos: a
          50\%, por exemplo, removem-se metade das cabeças e metade dos neurônios.
    \item \textbf{Uniforme:} cada camada recebe exatamente a mesma fração de poda,
          decidida isoladamente dentro dela. É a alocação mais simples possível, pois
          ignora qualquer diferença de importância entre camadas.
\end{itemize}

% AJUSTE: 24/10 (3) — pedido do Enzo: (i) o item Global explica por que as normas L1 não
% se comparam. Números calculados sobre os pesos pré-treinados públicos do GPT-2 small
% (openai-community/gpt2, model.safetensors), com as mesmas fórmulas de head_importances e
% mlp_importances (prune_structured.py): normas das 144 cabeças de 10.290,6 a 29.467,1
% (mediana 19.537,0); dos 36.864 neurônios, de 52,43 a 723,45 (mediana 143,60); menor
% cabeça / maior neurônio = 14,2; |w| médio por estrutura (mediana) 0,0993 na cabeça e
% 0,0934 no neurônio. Em ordenação conjunta, zero cabeças entre as 3.700, 11.102, 18.504 e
% 33.307 menores normas (10%, 30%, 50%, 90% de 37.008 estruturas). (ii) O "adianta-se"
% ganha o porquê: padrão das Tabelas tab:comp_magnitude e tab:comp_estruturada, e a
% interpretação de Resultados, subsec:perfis_divergem.
O contraste entre esses dois escopos é o que estabelece a linha de base contra a qual a
estratégia otimizada é julgada. Adianta-se que nenhum deles domina o outro em toda a faixa
de esparsidade, e é essa constatação que motiva a análise de sensibilidade. Na poda por
magnitude, o escopo uniforme obteve \textit{perplexity} menor entre 30\% e 70\%, e o
global, nas demais taxas; na poda estruturada, o global obteve \textit{perplexity} menor em
todas as taxas, exceto 20\% (Seções~\ref{sec:res_magnitude} e~\ref{sec:res_estruturada}).

Nenhum dos dois domina porque nenhum deles mede quanto cada camada tolera ser podada. O
escopo uniforme trata todas as camadas como igualmente tolerantes e poda a mesma fração de
cada uma, inclusive das mais sensíveis. O escopo global supõe que um peso, ou uma
estrutura, de menor magnitude é menos importante, qualquer que seja a camada em que se
encontra, e por isso concentra o corte nas camadas em que esses elementos predominam.
Quando essas camadas são também as que melhor toleram a poda, o escopo global leva
vantagem; quando não são, ele concentra a poda justamente em camadas sensíveis, e o
uniforme, que não a concentra em nenhuma, sai à frente. Como essa correspondência difere de
uma estratégia para outra e de uma taxa de poda para outra, a vantagem muda de lado ao
longo da faixa (Seção~\ref{subsec:perfis_divergem}). A análise de sensibilidade parte
dessa constatação para que a alocação deixe de depender de uma suposição: em vez de
inferir a tolerância de cada camada a partir da magnitude de seus pesos, ela a mede
diretamente, podando cada camada isoladamente e observando a degradação resultante
(Seção~\ref{subsec:fase1}).

\subsection{Taxas de Poda Avaliadas}\label{subsec:taxas}

% AJUSTE: 24/10 — seção reescrita para ficar mais didática, com os termos do restante do
% texto: "nível" → "taxa de poda"; "modelo base" → "modelo de referência"; "controle
% interno" passa a ser explicado; "ponto de colapso" (só aparecia aqui) vira a descrição
% do que se observa, com exemplo dos CSVs (magnitude global 26,87 a 20% e 165,20 a 40%;
% uniforme 104,82 a 40%; estruturada global 356,20 a 10% = 14x 25,17); "varredura", usada
% em todo o texto sem definição, é definida aqui; "recarregado do zero" (lia-se como
% retreinado) vira "a partir do modelo de referência original", como na Metodologia.
Cada estratégia de poda foi avaliada, nos dois escopos de alocação, global e uniforme
(Seção~\ref{subsec:escopos}), em dez taxas de poda: 0\%, 10\%, 20\% e assim por diante, até
90\%. Como definido na Seção~\ref{sec:metodologia_poda_sem}, a taxa de poda é a fração dos
elementos podáveis removida: pesos, na poda por magnitude, e cabeças e neurônios, na poda
estruturada. O conjunto das dez avaliações de uma mesma estratégia em um mesmo escopo é o
que se denomina, ao longo do texto, uma \textbf{varredura} --- por exemplo, a varredura
uniforme da poda por magnitude.

A taxa de 0\% corresponde ao modelo base, sem nenhuma poda. Ela foi incluída em
cada varredura como verificação: avaliado pelo mesmo procedimento das demais taxas, o
modelo base deve reproduzir a \textit{perplexity} medida na Etapa 1, 25,17, o que
ocorreu em todas as varreduras.

% AJUSTE: 24/10 (2) — pedido do Enzo: o parágrafo não antecipa resultados. Sai o exemplo
% numérico dos CSVs; a justificativa passa a ser a diferença de granularidade entre as
% estratégias (neurônio: 768 + 768 = 1.536 pesos; cabeça: 768x192 + 64x768 = 196.608).
% Termo padronizado em todo o texto: "modelo base" (antes "modelo de referência").
Avaliou-se toda a faixa, de 10\% a 90\%, porque não se sabia de antemão a partir de qual
taxa de poda cada estratégia passaria a degradar o modelo de forma acentuada. Não havia
razão para supor que essa taxa fosse a mesma nas duas estratégias, pois elas removem
elementos de tamanhos muito diferentes: a poda por magnitude remove pesos individuais,
enquanto a poda estruturada remove estruturas inteiras --- 1.536 pesos, no caso de um
neurônio do MLP, e 196.608, no caso de uma cabeça de atenção
(Seção~\ref{subsec:metodo_estruturado}). Avaliar apenas parte da faixa, como somente as
taxas baixas, poderia deixar de fora justamente a taxa em que a degradação de uma delas
se acentua.

Por fim, como estabelecido na Seção~\ref{sec:metodologia_poda_sem}, cada taxa de poda
partiu do modelo base original, e não do modelo já podado na taxa anterior: a
configuração de 30\%, por exemplo, foi obtida removendo 30\% dos elementos podáveis do
modelo base, e não acrescentando 20 pontos percentuais à poda do modelo de 10\%.
Assim, o resultado de uma taxa não depende das taxas avaliadas antes dela.

% AJUSTE: 12/10 — reunião com o orientador: a motivação da técnica (nenhum método da
% Etapa 2 mede a sensibilidade das camadas) e a descrição do procedimento passaram para a
% Metodologia (Seção sec:metodologia_poda_com, Figura fig:fluxo_sensibilidade). Aqui ficam
% a aplicação ao GPT-2: valores, forma exata da medida e regra de alocação.
\section{Etapa 3 --- Poda com Otimização}\label{sec:estrategia_otimizada}

Na terceira etapa, aplicou-se a técnica descrita na Seção~\ref{sec:metodologia_poda_com}
--- a análise de sensibilidade por camada, cujo procedimento está na
Figura~\ref{fig:fluxo_sensibilidade} --- às duas estratégias de poda, nas mesmas taxas de
poda da Etapa 2 (Seção~\ref{subsec:taxas}). As subseções trazem os valores adotados: a
medida da degradação e a taxa-sonda da fase 1, e a regra de alocação da fase 2, com o teto
e o parâmetro de intensidade.

\subsection{Fase 1 --- Perfil de Sensibilidade}\label{subsec:fase1}

O perfil foi levantado conforme o passo 1 do procedimento da
Seção~\ref{sec:metodologia_poda_com}: cada camada $\ell$ foi podada isoladamente, a uma
taxa-sonda fixa, e a \textit{perplexity} resultante foi medida. A degradação foi registrada
em escala logarítmica:

\begin{equation}\label{eq:delta_perfil}
\Delta_{\ell} = \ln\!\left(\frac{\text{PPL}_{\ell}}{\text{PPL}_{\text{base}}}\right),
\end{equation}

\noindent onde $\text{PPL}_{\ell}$ é a perplexity do modelo com apenas a camada $\ell$
podada. O logaritmo é necessário porque a perplexity de modelos podados sem recuperação
varia por ordens de grandeza; em escala linear, uma única camada catastrófica tornaria
todas as demais indistinguíveis entre si.

Adotou-se taxa-sonda de 50\%. O procedimento não envolveu treinamento e produziu, ao final, um
perfil com uma medida de sensibilidade por camada, para cada estratégia. Como as duas
estratégias removem coisas diferentes, os perfis foram levantados separadamente e não são
intercambiáveis --- ponto que a Seção~\ref{subsec:perfis_divergem} demonstra
experimentalmente.

% AJUSTE: 24/10 (3) — pedido do Enzo: "removem coisas diferentes" e "não são
% intercambiáveis" ganham a explicação. Sonda da magnitude: limiar único sobre as quatro
% matrizes da camada (prune_magnitude_per_layer, sensitivity.py). Conferido nos pesos
% pré-treinados: a 50%, cada cabeça perde de 34,6% a 69,3% de seus pesos e cada neurônio,
% de 25,0% a 91,3%; nenhum é anulado por inteiro. Sonda estruturada a 50%: int(0,5 x 12) = 6
% cabeças e int(0,5 x 3.072) = 1.536 neurônios (prune_structured_per_layer).
Os perfis não são intercambiáveis porque, à mesma taxa-sonda de 50\%, as duas estratégias
causam a uma camada danos de natureza diferente. A poda por magnitude anula metade dos
pesos podáveis da camada, escolhidos um a um entre os de menor valor absoluto, onde quer
que estejam. Os pesos anulados ficam distribuídos entre as cabeças e os neurônios, de modo
que nenhum deles é eliminado por inteiro: as 12 cabeças e os 3.072 neurônios continuam em
operação, cada um com parte de seus pesos. A poda estruturada, por sua vez, remove seis das
12 cabeças e 1.536 dos 3.072 neurônios, cada um com todos os seus pesos, pequenos ou
grandes, e mantém intactos os demais. A diferença se assemelha à de duas formas de cortar
pela metade as horas de trabalho de uma equipe: reduzir a jornada de todos os membros ou
dispensar metade deles. O total cortado é o mesmo, mas, no primeiro caso, todas as funções
da equipe continuam sendo exercidas, ainda que com menos tempo, e, no segundo, as funções
dos membros dispensados deixam de ser exercidas.

Em consequência, a degradação medida pelas duas estratégias reflete propriedades
diferentes da camada: na poda por magnitude, o quanto ela perde ao ter anulados seus pesos
de menor valor absoluto; na poda estruturada, o quanto ela perde ao ficar sem parte de suas
cabeças e de seus neurônios. Nada garante que uma camada que tolera bem um desses danos
tolere também o outro, e por isso o perfil de uma estratégia não serve para alocar a
outra: usá-lo seria distribuir a poda segundo a reação das camadas a um dano diferente do
que de fato se aplica. Cada perfil foi obtido podando as camadas com a própria estratégia,
para que medisse a reação ao mesmo dano que a alocação depois aplica, e orientou apenas a
alocação dessa estratégia.

% AJUSTE: 24/10 — pedido do Enzo: o parágrafo da segunda taxa-sonda era denso e não dizia
% por que ela é necessária nem por que é 30%. Desdobrado em três parágrafos: (i) a dúvida
% que motiva a réplica e por que ela importa; (ii) o raciocínio da comparação, com
% analogia; (iii) as condições que levam a 30%. O motivo original da escolha de 30% não
% estava registrado em nenhum lugar (texto, código, notas); a justificativa de (iii) foi
% formulada nesta revisão, sem recorrer a resultados. "Artefato do regime em que a sonda
% opera" sai por ser jargão.
A taxa-sonda de 50\%, porém, é uma escolha arbitrária: nenhum critério a torna preferível
a outras taxas próximas. Isso levanta uma dúvida: a ordenação das camadas --- quais são as
mais sensíveis e quais as mais robustas --- reflete uma característica das próprias
camadas ou apenas da taxa-sonda escolhida? A dúvida importa porque o perfil é medido uma
única vez, a 50\%, mas orienta a alocação em todas as taxas de poda avaliadas, de 10\% a
90\%. Se a ordenação mudasse com a taxa-sonda, uma camada classificada como robusta a
50\% poderia não sê-lo em outras taxas, e a alocação podaria mais justamente onde não
deveria.

Para responder a essa dúvida, o perfil completo foi medido novamente, nas duas
estratégias, a uma segunda taxa-sonda, e as duas ordenações foram comparadas. O raciocínio
é o mesmo de aplicar duas provas de dificuldades diferentes a uma mesma turma: se a
classificação dos alunos se repete nas duas, ela reflete os alunos, e não a prova. Do mesmo
modo, se as camadas aparecem na mesma ordem nas duas taxas-sonda, a ordenação é uma
propriedade das camadas.

A segunda taxa-sonda foi de 30\%, valor que atende a três condições. Primeiro, ela precisava
diferir o bastante de 50\% para que a comparação fosse informativa: com duas taxas muito
próximas, a coincidência das ordenações seria esperada e nada provaria. Segundo, convinha
que ficasse abaixo de 50\%, e não acima: acima dessa taxa, cada camada perde mais da
metade de seus elementos, e a medição se aproxima da remoção da camada inteira,
granularidade mais grossa que a adotada neste trabalho
(Seção~\ref{subsec:metodo_estruturado}). Terceiro, ela não podia ser baixa demais: a 10\%,
por exemplo, cada camada seria tão pouco alterada que as degradações ficariam todas
pequenas e próximas entre si, e a ordenação passaria a depender de diferenças mínimas. A
taxa de 30\%, 20 pontos percentuais abaixo da primeira, fica entre esses extremos.

A segunda medição serviu apenas como verificação e não foi usada em nenhuma alocação: todas
as taxas de poda por camada reportadas derivam do perfil de 50\%. O resultado da comparação
é apresentado na Seção~\ref{subsec:robustez_sonda}.

\subsection{Fase 2 --- Alocação Adaptativa de Taxas}\label{subsec:fase2}

% AJUSTE: 12/10 — reunião com o orientador: o fator α e o "orçamento global" não ficaram
% claros. "Orçamento" passa a "taxa de poda", e a regra é apresentada como generalização
% do exemplo da Metodologia (Tabela tab:exemplo_alocacao): o papel de α, do teto e de β
% é ilustrado com os mesmos números (Tabela tab:exemplo_beta).
% AJUSTE: 21/10 — acompanha a nova ordem das camadas do exemplo da Metodologia
% (robustezes 2, 5, 1 e 4): a camada mais robusta passa a ser a 2, e a mais sensível, a 3.
% AJUSTE: 24/10 — pedido do Enzo: seção reescrita para ser mais didática e menos técnica,
% só com termos que já aparecem no texto (fator, teto, parâmetro de intensidade, robustez,
% degradação, varredura uniforme). A regra passa a ser explicada em palavras, sobre o
% exemplo da Tabela tab:exemplo_alocacao, e a Equação eq:alocacao_dev fica no fim, como
% resumo (os Resultados a citam seis vezes). A média ponderada por w_l sai da equação:
% no GPT-2 ela coincide com a média simples, e o texto registra a ponderação em palavras.
% "Bisseção" vira "tentativas sucessivas" (o código, sensitivity.py:allocate_rates, faz
% bisseção). Piso da robustez: eps = 1e-4 no código; 0,0346 / 0,0001 = 346 vezes.
A fase 2 transforma o perfil de sensibilidade em uma alocação, isto é, em uma taxa de poda
para cada camada. Ela segue os passos 2 e 3 do procedimento da
Seção~\ref{sec:metodologia_poda_com}, já ilustrados pela Tabela~\ref{tab:exemplo_alocacao}
com um modelo hipotético de quatro camadas. Esta subseção retoma o mesmo exemplo para
apresentar a regra adotada e os três elementos que a compõem: o fator, o teto e o
parâmetro de intensidade.

A regra parte da ideia da Tabela~\ref{tab:exemplo_alocacao}: podar mais as camadas
robustas e menos as sensíveis. As robustezes indicam apenas a proporção entre as taxas por
camada. No exemplo, as robustezes 2, 5, 1 e 4 indicam que a camada 2 deve perder cinco
vezes mais elementos que a camada 3, mas não quanto cada uma perde. Esse valor é dado pelo
\textbf{fator}: todas as robustezes são multiplicadas por um mesmo número, escolhido para
que a média das taxas por camada seja igual à taxa de poda pretendida. À taxa de poda de
30\%, o fator é 0,10, e as taxas por camada são 20\%, 50\%, 10\% e 40\%, cuja média é
30\%. À taxa de 15\%, o fator cai pela metade, para 0,05, e todas as taxas também. Em
resumo, o perfil decide como a poda se distribui entre as camadas, e o fator decide quanto
se remove no total.

O \textbf{teto} limita a taxa de cada camada a 95\%. Sem ele, taxas de poda altas poderiam
levar a camada mais robusta a perder todos os seus elementos. No mesmo exemplo, à taxa de
poda de 60\%, o fator dobraria para 0,20, e a camada 2, de robustez 5, receberia 100\%:
seria inteiramente podada. Com o teto, ela fica em 95\%, e os 5 pontos percentuais que
deixa de perder passam às outras três camadas, para que a média continue em 60\%. O fator
sobe, então, para cerca de 0,207, e as camadas 1, 3 e 4 passam a 41,4\%, 20,7\% e 82,9\%.
Como, nessa situação, o fator não resulta de uma simples divisão, ele foi calculado por
tentativas sucessivas: ajusta-se o fator até que a média das taxas por camada coincida com
a taxa de poda.

O \textbf{parâmetro de intensidade}, representado pelo expoente $\beta$, controla o quanto
a alocação segue as diferenças do perfil. Com $\beta = 1$, as taxas seguem as robustezes na
proporção exata, como na Tabela~\ref{tab:exemplo_alocacao}. Com $\beta = 0{,}5$, cada
robustez é substituída por sua raiz quadrada, o que aproxima as taxas por camada entre si
sem alterar a ordem das camadas: a camada 2, que perdia cinco vezes mais elementos que a
camada 3, passa a perder 2,24 vezes mais (Tabela~\ref{tab:exemplo_beta}). Esse controle é
necessário porque não há limite para a diferença de robustez entre duas camadas: uma camada
que quase não degrada o modelo teria robustez muito alta e, sem amortecimento,
concentraria praticamente toda a poda. Os experimentos avaliaram $\beta = 1$ e
$\beta = 0{,}5$ sobre o mesmo perfil, o que permitiu comparar as duas intensidades sem
repetir a fase 1, a etapa de maior custo.

\begin{table}[htb]
\centering
\caption{Taxas de poda por camada no modelo hipotético da
Tabela~\ref{tab:exemplo_alocacao}, à taxa de poda de 30\%, para três valores do parâmetro
de intensidade $\beta$. Em todas as linhas a média é 30\%, e nenhuma camada atinge o
teto.}
\label{tab:exemplo_beta}
\small
\begin{tabular}{cc ccccc}
\hline
$\boldsymbol{\beta}$ & \textbf{Fator} & \textbf{Camada 1} & \textbf{Camada 2} &
\textbf{Camada 3} & \textbf{Camada 4} & \textbf{Média} \\
\hline
1,0 & 0,100 & 20,0\% & 50,0\% & 10,0\% & 40,0\% & 30\% \\
0,5 & 0,180 & 25,5\% & 40,3\% & 18,0\% & 36,1\% & 30\% \\
0   & 0,300 & 30,0\% & 30,0\% & 30,0\% & 30,0\% & 30\% \\
\hline
\end{tabular}
\end{table}

Para referência, a regra completa pode ser escrita de forma compacta na
Equação~\ref{eq:alocacao_dev}:

\begin{equation}\label{eq:alocacao_dev}
r_{\ell} = \min\!\left(\alpha \cdot \Delta_{\ell}^{-\beta},\ r_{\max}\right),
\qquad
\frac{1}{12}\sum_{\ell=0}^{11} r_{\ell} = s,
\end{equation}

\noindent em que $r_{\ell}$ é a taxa de poda da camada $\ell$; $\Delta_{\ell}$ é a sua
degradação (Equação~\ref{eq:delta_perfil}), cujo inverso é a robustez; $\alpha$ é o fator;
$\beta$, o parâmetro de intensidade; $r_{\max} = 0{,}95$, o teto; e $s$, a taxa de poda. A
segunda parte da equação exige que a média das taxas das 12 camadas seja igual à taxa de
poda. Na implementação, essa média é ponderada pelo número de parâmetros podáveis de cada
camada, o que permite aplicar a regra a modelos com camadas de tamanhos diferentes; no
GPT-2, cujas camadas têm o mesmo tamanho, ela coincide com a média simples.

O valor $\beta = 0$ também foi avaliado, mas como verificação, e não como alocação
candidata. Com ele, todas as robustezes passam a valer 1, todas as camadas recebem a mesma
taxa --- última linha da Tabela~\ref{tab:exemplo_beta} --- e a regra deve reproduzir a
varredura uniforme da Etapa 2. Se reproduzir, confirma-se que a regra está corretamente
implementada. Na poda estruturada, a reprodução deve ser exata. Na poda por magnitude, não:
a regra usa um único limiar para as quatro matrizes de cada camada, enquanto a varredura
uniforme usa um limiar para cada matriz. As duas removem a mesma fração de cada camada, mas
não exatamente os mesmos pesos. A verificação é apresentada na
Seção~\ref{subsec:controle_beta_zero}.

Por fim, a implementação inclui uma precaução que a equação não mostra. Se a poda de uma
camada na fase 1 não aumentasse a \textit{perplexity} --- o que pode ocorrer, pois há casos
em que remover cabeças de atenção a \textit{melhora} \cite{Michel2019} ---, a degradação
dessa camada seria nula ou negativa, e a robustez, que é o seu inverso, não poderia ser
calculada. Por isso, no cálculo da robustez, a degradação nunca fica abaixo de um valor
mínimo positivo, 0,0001, e uma camada nessa situação entra na alocação como a mais robusta.
Nos perfis usados nas alocações, essa precaução não chegou a ser acionada: a menor
degradação foi de 0,0346, mais de 300 vezes esse valor mínimo.

% AJUSTE: revisão 30/09 — "fecha essa lacuna", marcado pelo orientador na Metodologia,
% aparecia também aqui; reescrito do mesmo modo.
% AJUSTE: 05/10 — revisão 30/09, "não consegui entender exatamente como é feito o
% fine-tuning de recuperação". A Metodologia remete a esta seção o protocolo concreto
% (configurações, número de passos, hiperparâmetros), que não estava aqui: os valores
% constavam só dos Resultados, e o texto dizia "para cada configuração", embora apenas
% sete tenham sido retreinadas. A seção passa a descrever o que codigo/finetune.py faz e
% o que codigo/resultados/finetune.csv registra. Saem duas afirmações que o código não
% sustenta: que o acúmulo de gradiente e a precisão mista foram necessários para caber
% na memória da GPU.
\section{Etapa 4 --- \textit{Fine-tuning} de Recuperação}\label{sec:finetuning}

As Etapas 2 e 3 operaram em regime \textit{one-shot}, isto é, podaram e avaliaram sem
retreinamento, de modo a isolar o dano causado pela poda. Esse regime, porém, não
corresponde ao procedimento da literatura: tanto \citeonline{Han2015} quanto
\citeonline{Li2017} retreinam o modelo após podá-lo. Na quarta etapa, esse retreinamento
foi incorporado ao experimento: sete configurações foram submetidas ao \textit{fine-tuning}
de recuperação, e a \textit{perplexity} de cada uma foi medida antes e depois do
retreinamento, o que indica quanto da degradação é recuperável. As subseções a seguir
apresentam as configurações retreinadas, o procedimento aplicado a cada uma, os cuidados
que condicionam a validade da medição e os parâmetros do treinamento.

\subsection{Configurações Retreinadas}\label{subsec:ft_configuracoes}

O retreinamento foi aplicado às seguintes configurações das Etapas 2 e 3:

\begin{itemize}
    \item \textbf{Controle denso:} o modelo base, sem poda, cuja função é
          descrita na Seção~\ref{subsec:ft_validade};
    \item \textbf{Poda não estruturada --- poda por magnitude}, com escopo uniforme, a
          30\% e a 50\% de esparsidade;
    \item \textbf{Poda estruturada --- poda de cabeças de atenção e de neurônios do MLP},
          com os dois alvos simultâneos e escopo global, a 10\% e a 30\% de estruturas
          removidas;
    \item \textbf{Poda estruturada com análise de sensibilidade}, com a alocação definida
          pelo perfil da própria poda estruturada, com $\beta = 1$, nas mesmas duas taxas.
\end{itemize}

% AJUSTE: 24/10 — pedido do Enzo: o parágrafo da escolha das configurações era um bloco
% único e denso. Passa a responder, em ordem, a três perguntas (qual escopo, em qual
% estratégia retreinar a alocação com análise de sensibilidade, em quais taxas), com o
% motivo de cada escolha em palavras simples. "Faixa em que a escolha de escopo tem
% consequência prática" vira a explicação do que ocorre fora dela (Resultados: até 20% a
% poda é quase gratuita; a 80-90% ambos os modelos estão destruídos); "nível" → "taxa";
% a razão da estruturada (única que reduz custo) ganha o porquê da magnitude não reduzir
% (pesos zerados, não removidos, Seção 5.3.3).
As configurações da lista foram escolhidas a partir dos resultados das Etapas 2 e 3 em
regime \textit{one-shot}, apresentados no Capítulo~\ref{chap:resultados}. A escolha
respondeu a três perguntas.

A primeira é qual dos dois escopos sem análise de sensibilidade, global ou uniforme,
retreinar em cada estratégia. Escolheu-se, em cada uma, o de menor \textit{perplexity} nas
taxas de poda em que a escolha faz diferença. Na poda por magnitude, foi o uniforme, que
supera o global entre 30\% e 70\%. Fora dessa faixa, a escolha pouco importa: até 20\%, os
dois quase não degradam o modelo, e a partir de 80\% ambos já o destroem
(Seção~\ref{sec:res_magnitude}). Na poda estruturada, foi o global, que supera o uniforme
em todas as taxas, exceto 20\% (Seção~\ref{sec:res_estruturada}).

A segunda é em qual estratégia retreinar também a alocação com análise de sensibilidade.
Escolheu-se apenas a poda estruturada, por ter sido a única a reduzir o tempo de inferência
e o consumo de energia (Seção~\ref{sec:res_custo}). Na poda por magnitude, os pesos são
zerados, mas não removidos (Seção~\ref{subsec:metodo_magnitude}), e o custo não cai. Usou-se
o parâmetro de intensidade $\beta = 1$, o de melhor desempenho na poda estruturada
(Seção~\ref{subsec:ablacao_beta}).

A terceira é em quais taxas de poda retreinar. Na poda estruturada, escolheram-se 10\% e
30\%. São as mesmas taxas nas duas alocações, global e com análise de sensibilidade, o que
permite compará-las também depois do retreinamento, e ambas estão na faixa de 10\% a 50\%
em que, em regime \textit{one-shot}, a alocação com análise de sensibilidade obteve
\textit{perplexity} menor que a global (Seção~\ref{sec:res_estruturada}). Na poda por
magnitude, escolheram-se duas taxas com danos de intensidades diferentes: 30\%, a primeira
taxa acima da região utilizável em regime \textit{one-shot} --- as taxas de até 20\%, em que
a poda quase não degrada o modelo ---, e 50\%, em que a degradação já é severa, com
\textit{perplexity} de 473,28 (Seção~\ref{sec:res_magnitude}). Com isso, o retreinamento é
observado tanto sobre um dano moderado quanto sobre um dano severo.

\subsection{Procedimento}\label{subsec:ft_procedimento}

Cada configuração passou pelo mesmo procedimento, implementado no módulo
\texttt{finetune} (Tabela~\ref{tab:modulos}):

\begin{enumerate}
    \item o modelo base foi recarregado a partir dos pesos pré-treinados e podado
          pelas mesmas funções das Etapas 2 e 3, de modo que o modelo retreinado fosse
          idêntico ao avaliado em regime \textit{one-shot};
    \item mediu-se a \textit{perplexity} do modelo podado, ponto de partida da
          recuperação; nas sete configurações, o valor obtido reproduziu, em todos os
          dígitos registrados, o das varreduras \textit{one-shot};
    \item na poda não estruturada, capturou-se a máscara binária dos pesos
          remanescentes (Seção~\ref{subsec:ft_validade});
    \item o modelo foi retreinado sobre a partição de treino do WikiText-2, com os
          parâmetros da Tabela~\ref{tab:finetune_protocolo};
    \item o modelo retreinado foi avaliado na partição de teste por todas as métricas da
          Etapa 1 (Tabela~\ref{tab:metricas}), pelo mesmo procedimento aplicado aos
          modelos não retreinados.
\end{enumerate}

% AJUSTE: 24/10 (3) — pedido do Enzo: o porquê de retreinar todos os parâmetros. Conferido
% em finetune.py: nenhum parâmetro é congelado; o AdamW recebe todos (_param_groups), e
% vieses e normalização só ficam sem weight decay. Empilhamento: subsec:bloco_transformer;
% pesos compartilhados entre lm_head e wte e papel da normalização: parametros_podaveis.
O retreinamento atualizou todos os parâmetros do modelo podado --- inclusive
\textit{embeddings}, vieses e parâmetros de normalização ---, e não apenas as matrizes
podáveis (Seção~\ref{subsec:parametros_podaveis}), já que o efeito da poda não fica
restrito às matrizes de onde os elementos foram removidos. As camadas do GPT-2 são
empilhadas (Seção~\ref{subsec:bloco_transformer}): a saída de cada uma é a entrada da
seguinte, e a saída da última alimenta a normalização final e a cabeça de linguagem. Ao
podar uma camada, altera-se a representação que ela entrega, e todos os componentes
posteriores, podados ou não, passam a receber entradas diferentes daquelas para as quais
foram pré-treinados. Recuperar o modelo exige, portanto, não só reajustar os pesos que
restaram nas matrizes podadas, mas também readaptar esses componentes às novas entradas.
Os parâmetros de normalização, por exemplo, controlam a escala de cada dimensão da entrada
da atenção e do MLP e podem reajustá-la às representações alteradas; restringir o
retreinamento às matrizes podáveis os manteria fixos nos valores do pré-treinamento,
calibrados para o modelo sem poda. Do mesmo modo, a cabeça de linguagem, que no GPT-2
compartilha os pesos com o \textit{embedding} de \textit{token}, precisa readaptar-se às
representações que recebe da última camada, o que só é possível se esse
\textit{embedding} também for atualizado. Atualizar todos os parâmetros, por outro lado, não desfaz a poda: na poda
estruturada, as estruturas removidas deixaram de existir, e, na poda por magnitude, os
pesos anulados foram mantidos em zero por uma máscara (Seção~\ref{subsec:ft_validade}).

O modelo avaliado é o obtido ao final
do último passo do otimizador: a partição de validação não foi utilizada, e não houve
parada antecipada nem seleção do melhor ponto do treinamento. Com isso, a quantidade de
retreinamento é idêntica em todas as configurações, condição para que suas
\textit{perplexities} pós-retreinamento sejam comparáveis entre si.

\subsection{Controle Denso e Preservação da Esparsidade}\label{subsec:ft_validade}

Duas decisões de implementação merecem registro por afetarem a validade dos resultados.

A primeira é o \textbf{controle denso}. O GPT-2 não foi pré-treinado no WikiText-2, de
modo que o simples \textit{fine-tuning} do modelo base no WikiText-2 já reduz a perplexity
por adaptação de domínio, sem poda alguma envolvida. Sem essa linha de controle, seria
impossível separar a recuperação do dano da poda da mera adaptação à base de dados. O modelo
base foi, portanto, retreinado pelo mesmo número de passos e avaliado junto às demais
configurações.

A segunda é a \textbf{preservação da esparsidade}. Na poda não estruturada os pesos são
apenas zerados, e o gradiente os restauraria já no primeiro passo do otimizador; ao final
do treinamento, o modelo estaria novamente denso. Como em \citeonline{Han2015}, manteve-se
uma máscara binária fixa, capturada logo após a poda e reaplicada depois de cada passo do
otimizador. Registre-se que zerar o gradiente não seria suficiente: o otimizador utilizado
possui momento e regularização por decaimento de pesos, e ambos alteram um parâmetro mesmo
quando seu gradiente é nulo. A esparsidade medida ao final do treinamento foi registrada
como verificação de integridade e coincidiu, nas duas configurações de magnitude, com a
medida logo após a poda (Seção~\ref{subsec:ft_verificacoes}). Na poda estruturada não
houve máscara, pois as estruturas já haviam sido fisicamente removidas.

\subsection{Parâmetros do Treinamento}\label{subsec:ft_protocolo}

Para o treinamento, as linhas da partição de treino do WikiText-2 foram concatenadas em
um único texto, tokenizado e dividido em blocos contíguos e disjuntos de 512
\textit{tokens}; o resto final, menor que um bloco, foi descartado. Resultaram 4.743
blocos, ou 2.428.416 \textit{tokens}. A divisão difere da adotada na avaliação, que
percorre o texto por janelas deslizantes sobrepostas (Seção~\ref{sec:metricas}): na
avaliação, a sobreposição garante contexto anterior a cada \textit{token} avaliado; no
treinamento, blocos disjuntos fazem de cada bloco uma amostra independente. A
Tabela~\ref{tab:finetune_protocolo} reúne os parâmetros adotados, idênticos nas sete
configurações.

\begin{table}[htb]
\centering
\caption{Parâmetros do \textit{fine-tuning} de recuperação.}
\label{tab:finetune_protocolo}
\small
\begin{tabular}{p{7.5cm} r}
\hline
\textbf{Parâmetro} & \textbf{Valor} \\
\hline
Blocos de treino (512 \textit{tokens} cada) & 4.743 \\
Blocos por lote & 4 \\
Lotes acumulados por passo do otimizador & 4 \\
Lote efetivo por passo & 16 blocos (8.192 \textit{tokens}) \\
Passos do otimizador & 297 (uma época) \\
Otimizador & AdamW \\
Taxa de aprendizado máxima & $5 \times 10^{-5}$ \\
Aquecimento linear da taxa de aprendizado & 17 passos (6\% do total, truncado) \\
Decaimento da taxa de aprendizado & linear, até zero no passo 297 \\
Decaimento de pesos & 0,01 \\
Limite da norma do gradiente & 1,0 \\
Precisão numérica & mista (FP16 e FP32) \\
Semente aleatória & 42 \\
\hline
\end{tabular}\\
\textit{Fonte: Elaboração própria, com valores de \texttt{codigo/finetune.py} e de
\texttt{codigo/resultados/finetune.csv}.}
\end{table}

Os gradientes de quatro lotes consecutivos, de quatro blocos cada, foram acumulados antes
de cada passo do otimizador, o que equivale a um lote efetivo de 16 blocos; a perda de
cada lote foi dividida por quatro, de modo que o gradiente acumulado correspondesse à
média dos lotes, e não à soma. O número de passos foi calculado para corresponder a uma
época, isto é, a uma passagem pelo conjunto de treino: os 4.743 blocos, divididos em lotes efetivos de 16,
resultam em 297 passos, com arredondamento para cima.

O otimizador AdamW \cite{Loshchilov2019} aplica o decaimento de pesos separadamente da
atualização pelo gradiente. O decaimento foi aplicado apenas às matrizes de peso,
deixando de fora vieses e parâmetros de normalização, conforme prática consolidada. A
taxa de aprendizado cresceu linearmente durante os 17 primeiros passos e decresceu
linearmente até zero no último, e a norma do gradiente foi limitada a 1,0 antes de cada
passo. A taxa de aprendizado e o lote efetivo estão entre os valores recomendados por
\citeonline{Devlin2019} para o \textit{fine-tuning} de \textit{Transformers}
pré-treinados --- taxa de $2 \times 10^{-5}$ a $5 \times 10^{-5}$ e lotes de 16 ou 32
exemplos ---; o decaimento de pesos de 0,01 e o aquecimento seguido de decaimento linear
são os adotados por esses autores no pré-treinamento do BERT. Nenhum desses valores foi
ajustado para este experimento (Seção~\ref{sec:limitacoes}). A quantidade de
retreinamento, uma época, fica abaixo das duas a quatro épocas recomendadas pelos mesmos
autores; o que essa limitação impede de concluir é discutido na
Seção~\ref{subsec:ft_verificacoes}.

A passagem direta e a retropropagação executaram em \textbf{precisão mista}
\cite{Micikevicius2018}: as operações foram realizadas em meia precisão (FP16), com os
pesos mantidos em precisão simples (FP32) e escalonamento dinâmico da perda, que evita o
subtransbordamento dos gradientes em meia precisão. A escolha aproveita as unidades de
cálculo especializadas da GPU Tesla T4, que alcança 65 TFLOPS em precisão mista contra
8,1 TFLOPS em precisão simples \cite{NVIDIA2019T4}, e corresponde à configuração
pressuposta na estimativa de memória do critério C3
(Seção~\ref{subsec:metodologia_selecao_modelo}). Ela introduz, porém, uma fonte adicional
de variação numérica que não existe nas etapas \textit{one-shot}, todas executadas em
precisão simples. O consumo de energia do treinamento foi medido pela biblioteca
CodeCarbon separadamente do consumo da avaliação: trata-se de um custo pago uma única vez,
que precisa ser amortizado contra a economia por inferência para que o balanço energético
do trabalho feche corretamente.

Diferentemente de todas as etapas anteriores, esta é estocástica --- a ordem dos lotes e o
\textit{dropout} introduzem variabilidade --- e, portanto, é a única em que a semente
aleatória efetivamente altera o resultado. Cada configuração foi retreinada uma única vez,
com a semente 42, reiniciada antes de cada configuração. A ordem dos lotes, sorteada por um
gerador próprio iniciado com essa semente, foi a mesma nas sete configurações, de modo que
as diferenças entre elas não decorrem da ordem em que os dados foram apresentados. A
ausência de repetições com sementes distintas é discutida entre as limitações
(Seção~\ref{sec:limitacoes}).

\section{\textit{Framework} de Implementação}\label{sec:framework}

A implementação foi feita em Python, com o \textit{framework} PyTorch \cite{Paszke2019}
para construção, manipulação e treinamento dos modelos. O PyTorch foi escolhido por seu
sistema de diferenciação automática, que constrói dinamicamente o grafo computacional e
simplifica a implementação do \textit{backpropagation} durante o \textit{fine-tuning}, e
por seu suporte nativo à aceleração em GPU.

% AJUSTE: 05/10 — parágrafo trazido da Metodologia (Etapa 1), segundo o esboço do
% orientador: as ferramentas usadas pertencem ao Desenvolvimento.
Os pesos pré-treinados foram obtidos na plataforma Hugging Face por meio da biblioteca
\texttt{transformers} \cite{Wolf2020}, que implementa sobre o PyTorch as arquiteturas dos
modelos candidatos. A Hugging Face não é uma biblioteca, mas um ambiente colaborativo de
hospedagem e distribuição de modelos, que hospedava mais de 3,1 milhões de modelos em 28
set. 2026, entre eles diversas famílias de modelos de linguagem, como GPT-2, OPT, Pythia,
Llama e BERT; a biblioteca que executa os modelos é o PyTorch.

As demais bibliotecas empregadas são:

\begin{itemize}
    \item \texttt{datasets}: carregamento da base de dados WikiText a partir da plataforma
          Hugging Face;
    \item \texttt{thop}: contagem de operações de ponto flutuante por análise do grafo
          computacional;
    \item \texttt{CodeCarbon} \cite{CodeCarbon2023}: medição empírica do consumo de
          energia e estimativa das emissões de $CO_2$, a partir do monitoramento do
          \textit{hardware} combinado à intensidade de carbono da rede elétrica local.
\end{itemize}

Uma observação relevante sobre a implementação da poda: ambas as estratégias foram
implementadas manualmente, sem recurso a bibliotecas dedicadas. No caso da poda
estruturada, essa decisão foi imposta pelo ambiente --- as versões mais recentes da
biblioteca \texttt{transformers} removeram a interface de poda de cabeças que existia em
versões anteriores. A remoção física das estruturas foi, portanto, realizada diretamente:
as matrizes foram reconstruídas com as dimensões reduzidas e a contabilidade interna do
módulo de atenção foi atualizada. A implementação da poda é coberta por uma
suíte de testes automatizados (\texttt{test\_poda}), executada sobre um GPT-2 de dimensões
reduzidas e pesos aleatórios --- o que está sob verificação é a mecânica da poda, não os
pesos pré-treinados. O teste central é o de \textbf{equivalência numérica}: remover
fisicamente um conjunto de cabeças de atenção, ou de neurônios do MLP, deve produzir
exatamente as mesmas saídas que apenas anular os pesos correspondentes no modelo de
dimensões originais. A equivalência é matematicamente necessária --- anular as fatias de
consulta, chave e valor de uma cabeça zera seu vetor de valores, de modo que sua
contribuição à saída da atenção se anula qualquer que seja a distribuição de atenção
calculada ---, e por isso qualquer erro de indexação na reconstrução das matrizes a
quebraria. Os demais testes verificam que a contabilidade interna do módulo de atenção
permanece coerente com as formas dos tensores, que a salvaguarda preserva ao menos uma
estrutura por camada mesmo sob taxa integral, que as estruturas removidas são efetivamente
as de menor norma $L_1$, que a poda por magnitude atinge a esparsidade solicitada
preservando as dimensões, que a regra de alocação respeita a taxa de poda e poda menos a
camada mais sensível, e que a máscara de esparsidade sobrevive a um passo do otimizador.

\section{\textit{Pipeline} Experimental}\label{sec:pipeline}

O código está organizado em módulos de responsabilidade única, o que permite que cada
estratégia de poda reutilize exatamente a mesma rotina de avaliação --- condição para que
as comparações sejam legítimas. A Tabela~\ref{tab:modulos} descreve os módulos e a
Figura~\ref{fig:pipeline_experimental} representa o fluxo experimental.

\begin{table}[htb]
\centering
\caption{Organização modular da implementação.}
\label{tab:modulos}
\small
\begin{tabular}{p{3.4cm} p{9.6cm}}
\hline
\textbf{Módulo} & \textbf{Responsabilidade} \\
\hline
\texttt{config} & Configuração central: modelo, base de dados, parâmetros de avaliação, semente e ambiente. \\
\texttt{utils} & Fixação das sementes e escrita incremental dos resultados em arquivo. \\
\texttt{data} & Carregamento e tokenização da base de dados. \\
\texttt{model} & Carregamento do GPT-2 pré-treinado. \\
\texttt{eval} & Perplexity por janela deslizante e métricas de custo computacional. \\
\texttt{energy} & Medição de energia e emissões. \\
\texttt{prune\_magnitude} & Poda não estruturada por magnitude, nos dois escopos. \\
\texttt{prune\_structured} & Poda estruturada de cabeças e neurônios, com remoção física. \\
\texttt{sensitivity} & Perfil de sensibilidade (fase 1) e alocação adaptativa (fase 2). \\
\texttt{finetune} & \textit{Fine-tuning} de recuperação, com preservação de esparsidade. \\
\texttt{run} & Orquestração da avaliação do modelo base. \\
\texttt{test\_poda} & Testes de corretude da poda, descritos na Seção~\ref{sec:framework}. \\
\hline
\end{tabular}\\
\textit{Fonte: Elaboração própria.}
\end{table}

\begin{figure}[htb]
\centering
\caption{Fluxo experimental do trabalho. Todas as estratégias convergem para a mesma
rotina de avaliação, o que garante a simetria necessária às comparações.}
\label{fig:pipeline_experimental}
\begin{tikzpicture}[font=\scriptsize, node distance=0.55cm and 0.45cm,
  fonte/.style={rectangle,rounded corners=2pt,draw,thick,fill=gray!8,
                minimum width=3.1cm,minimum height=0.75cm,align=center},
  etapa/.style={rectangle,rounded corners=2pt,draw,thick,fill=gray!18,
                minimum width=3.3cm,minimum height=0.9cm,align=center},
  saida/.style={rectangle,rounded corners=2pt,draw,thick,fill=gray!30,
                minimum width=4.6cm,minimum height=0.75cm,align=center}]

  % --- entradas ---
  \node[fonte] (mod) at (-2.1,0) {GPT-2 \textit{small}\\ \scriptsize pré-treinado};
  \node[fonte] (dat) at (2.1,0) {WikiText-2\\ \scriptsize base de dados de avaliação};

  % --- baseline ---
  \node[etapa] (base) at (0,-1.35) {\textbf{Modelo base}\\ \scriptsize perplexity 25,17};
  \draw[-{Stealth}] (mod.south) -- ++(0,-0.3) -| (base.north);
  \draw[-{Stealth}] (dat.south) -- ++(0,-0.3) -| (base.north);

  % --- três estratégias ---
  \node[etapa] (mag) at (-4.6,-3.0) {\textbf{Magnitude}\\ \scriptsize global / uniforme};
  \node[etapa] (est) at (0,-3.0)    {\textbf{Estruturada}\\ \scriptsize global / uniforme};
  \node[etapa] (sen) at (4.6,-3.0)  {\textbf{Sensibilidade}\\ \scriptsize perfil $\rightarrow$ alocação};
  \foreach \n in {mag,est,sen} \draw[-{Stealth}] (base.south) -- ++(0,-0.35) -| (\n.north);

  % --- fine-tuning ---
  \node[etapa] (ft) at (0,-4.75) {\textbf{\textit{Fine-tuning} de recuperação}\\
        \scriptsize com controle denso e máscara de esparsidade};
  \foreach \n in {mag,est,sen} \draw[-{Stealth}] (\n.south) -- ++(0,-0.35) -| (ft.north);

  % --- avaliação ---
  \node[etapa] (av) at (0,-6.3) {\textbf{Avaliação comum}\\
        \scriptsize perplexity, parâmetros, tamanho, FLOPs, tempo, energia};
  \draw[-{Stealth}] (ft) -- (av);

  \node[saida] (out) at (0,-7.6) {Tabelas e figuras do Capítulo~\ref{chap:resultados}};
  \draw[-{Stealth}] (av) -- (out);

  % --- desvio one-shot ---
  \draw[-{Stealth},dashed,gray!70] (mag.west) -- ++(-0.55,0)
     |- node[left=2pt,align=right,font=\scriptsize,black,pos=0.25]
        {regime\\ \textit{one-shot}} (av.west);
\end{tikzpicture}
\fonte{Elaborada pelo autor.}
\end{figure}

Como indica a linha tracejada da figura, as varreduras \textit{one-shot} contornam a etapa
de \textit{fine-tuning} e vão direto à avaliação. Essa bifurcação é o que permite separar,
nos resultados, o dano causado pela poda daquilo que o retreinamento consegue recuperar.

\section{Ambiente de Execução e Reprodutibilidade}\label{sec:ambiente}

Os experimentos foram executados na plataforma Kaggle, em sessões com GPU NVIDIA Tesla T4. A
implementação detecta automaticamente o dispositivo disponível, o que permite executar
localmente em CPU para testes rápidos sem alteração de código.

Quanto à reprodutibilidade, três aspectos merecem registro.

O primeiro é que as varreduras \textit{one-shot} são \textbf{determinísticas}. Não há
treinamento, o modelo opera em modo de avaliação, o conjunto de teste é fixo e os critérios
de poda são ordenações determinísticas sobre pesos fixos. Verificou-se empiricamente que
execuções independentes, em sessões distintas, reproduziram a perplexity até a última casa
decimal. Repetições com sementes diferentes, nessas etapas, produziriam linhas idênticas e
não acrescentariam informação. A semente passa a importar apenas no \textit{fine-tuning}.

O segundo é que as medições de \textbf{tempo e energia dependem da sessão}. Observou-se que
o mesmo modelo base, sem qualquer poda, registrou tempos e consumos sensivelmente
distintos em sessões diferentes da plataforma, o que se atribui à variação de
\textit{hardware} físico e de carga concorrente. Por essa razão, comparações de custo entre
varreduras executadas em sessões distintas foram sempre feitas em termos relativos ao modelo
base da própria sessão, nunca em valores absolutos.

O terceiro é que cada configuração experimental gravou uma linha em arquivo de resultados,
contendo não apenas as métricas mas também os parâmetros que a geraram --- estratégia,
escopo, taxa alvo, taxa efetiva e, no caso da estratégia otimizada, as taxas atribuídas a
cada camada. Isso permite rastrear qualquer número apresentado no
Capítulo~\ref{chap:resultados} até a configuração exata que o produziu.

\section{Limitações Metodológicas}\label{sec:limitacoes}

Reconhecem-se, desde o desenho experimental, limitações que condicionam o escopo dos
resultados:

\begin{itemize}
    \item \textbf{Um único modelo e uma única base de dados.} Os experimentos empregaram
          exclusivamente o GPT-2 \textit{small} sobre o WikiText-2. Os resultados podem
          não se transferir diretamente para modelos de escala substancialmente maior ou
          para bases de dados de natureza distinta. A implementação foi mantida parametrizável
          justamente para viabilizar essa verificação em trabalhos futuros.

    \item \textbf{Duas estratégias de poda, e não mais.} A restrição de tempo experimental
          levou à escolha de duas estratégias que cobrem os extremos de granularidade.
          Métodos baseados em informação de gradiente ou de segunda ordem ficaram fora do
          escopo.

    \item \textbf{Faixa restrita de taxas-sonda no perfil de sensibilidade.} O perfil que
          alimenta a alocação foi levantado a uma única taxa, e sua estabilidade foi
          verificada por meio de uma réplica a uma segunda taxa
          (Seção~\ref{subsec:robustez_sonda}). Ambas pertencem, porém, ao mesmo regime
          intermediário de poda: nada garante que a ordenação de sensibilidade entre
          camadas se mantenha em regimes muito mais brandos ou muito mais agressivos.

    \item \textbf{Sensibilidade medida em isolamento.} A sonda mede o efeito de cada camada
          isoladamente e ignora interações entre elas. Como se verifica no
          Capítulo~\ref{chap:resultados}, a degradação é fortemente superaditiva, de modo
          que o perfil deve ser lido como ordenação relativa de importância, e não como
          modelo preditivo do desempenho de uma configuração completa.

    \item \textbf{Dependência do \textit{hardware}.} Tempo e energia são intrinsecamente
          dependentes da máquina. Os valores absolutos não devem ser comparados entre
          ambientes distintos; a análise apoia-se em comparações relativas dentro de um
          mesmo ambiente.

    % AJUSTE: 05/10 — "e de experimentos preliminares" sai: o código e o histórico não
    % registram nenhum; os valores padrão de finetune.py nunca mudaram desde o commit
    % f95cf76, e a única execução anterior foi o teste de fumaça (2 passos, descartado).
    \item \textbf{Hiperparâmetros não otimizados.} Os hiperparâmetros do
          \textit{fine-tuning} foram tomados de valores recomendados na literatura
          (Seção~\ref{subsec:ft_protocolo}), sem busca sistemática. Valores melhores
          podem existir.

    \item \textbf{Semente única no \textit{fine-tuning}.} Esta é a única etapa estocástica
          do experimento, e todas as configurações foram executadas com uma única semente.
          As perplexities pós-retreinamento não vêm, portanto, acompanhadas de estimativa
          de variância, e diferenças pequenas entre configurações retreinadas devem ser
          lidas com a cautela correspondente. A comparação principal da etapa --- alocação
          com análise de sensibilidade contra alocação global --- é afetada por essa
          limitação justamente na taxa de poda em que as duas praticamente empatam.

    \item \textbf{Estimativa de emissões.} A conversão de energia em emissões depende de
          tabelas regionais de intensidade de carbono, que carregam incerteza e variaram
          entre as sessões de execução. As emissões devem ser lidas como aproximações de
          ordem de grandeza.
\end{itemize}

Apesar dessas limitações, o desenho descrito é suficiente para responder às duas perguntas
enunciadas na Seção~\ref{sec:visao_geral}, que são o objeto do
Capítulo~\ref{chap:resultados}.
